\documentclass{report}
\usepackage[T1]{fontenc}
\usepackage[utf8]{inputenc}
\usepackage[backend=biber, style=ieee]{biblatex}
\usepackage{csquotes}
\usepackage[portuguese]{babel}
\usepackage{blindtext}
\usepackage[printonlyused]{acronym}
\usepackage{hyperref}
\usepackage{graphicx}
\usepackage{indentfirst}
\usepackage{comment}
\usepackage{booktabs}
\usepackage{tabularx}
\usepackage{longtable}
\usepackage{multirow} 
\usepackage[table]{xcolor}
\usepackage[
    a4paper,        % Define o papel como A4
    left=3cm,       % Margem esquerda
    right=3cm,      % Margem direita
    top=1.5cm,      % Margem superior
    bottom=1.5cm,   % Margem inferior
]{geometry}

\setcounter{tocdepth}{1}
\bibliography{bibliografia}


\begin{document}
%%
% Definições
%
\def\titulo{Google Assistant - IM Report}
\def\data{DATA}
\def\autores{Duarte Santos, Henrique Teixeira}
\def\autorescontactos{(113304) duartesantos@ua.pt, (114588) henriqueft04@ua.pt}
\def\versao{VERSAO}
\def\departamento{Dept. de Eletrónica, Telecomunicações e Informática}
\def\empresa{Universidade de Aveiro}
\def\logotipo{./Images/ua.pdf}
%
%%%%%% CAPA %%%%%%
%
\begin{titlepage}

\begin{center}
%
\vspace*{50mm}
%
{\Huge\textbf{\titulo}} \\
\vspace*{10mm}
{\Large \departamento\\ \empresa}
%
\vspace{10mm}
%
{\Large \empresa}\\
%

\vspace{20mm}
%
\begin{figure}[h]
\center
\includegraphics{\logotipo}
\end{figure}

\vspace{30mm}
%
{\LARGE \autores}\\ 
{\Large \autorescontactos}\\ 
%
%
\vspace{30mm}
\end{center}
%

\end{titlepage}



\pagenumbering{roman}

%%%%%% RESUMO %%%%%%
%\begin{abstract}
%Resumo de 200-300 palavras.
%\end{abstract}

%%%%%% Agradecimentos %%%%%%
% Segundo glisc deveria aparecer após conclusão...
%\renewcommand{\abstractname}{Agradecimentos}
%\begin{abstract}
%Eventuais agradecimentos.
%Comentar bloco caso não existam agradecimentos a fazer.
%\end{abstract}


\tableofcontents
% \listoftables     % descomentar se necessário
% \listoffigures    % descomentar se necessário


%%%%%%%%%%%%%%%%%%%%%%%%%%%%%%%
\clearpage
\pagenumbering{arabic}

%%%%%%%%%%%%%%%%%%%%%%%%%%%%%%%%
\chapter{Introdução}
\label{chap.introducao}

O sistema desenvolvido consiste num assistente multimodal que controla o Google Maps por voz e/ou por gestos por um utilizador. De forma a garantir uma interação fluida, o assistente fornece feedback auditivo a cada ação executada pelo utilizador. O projeto foi implementado na linguagem Python, com recurso à biblioteca Selenium. A escolha desta tecnologia deve-se à capacidade do Selenium automatizar e controlar o browsing na web.

Para os componentes de interação (voz e gestos) e processamento central (Motor de Fusão), foi adotada a framework multimodal do IEETA. \\

O projeto foi desenvolvido em três iterações distintas:

\begin{itemize}
    \item \textbf{Iteração 1 - Voz:} Implementação de um assistente controlado por voz utilizando RASA NLU para reconhecimento de intenções;

    \item \textbf{Iteração 2 - Gestos:} Adição de reconhecimento de gestos através de Kinect e GenericGesturesModality;

    \item \textbf{Iteração 3 - Fusão:} Integração das duas modalidades (voz e gestos) através de um motor de fusão baseado em SCXML. Esta iteração permite ao utilizador interagir com o Google Maps usando voz, gestos ou ambos simultaneamente.
\end{itemize}


%%%%%%%%%%%%%%%%%%%%%%%%%%%%%%%%
\chapter{Arquitetura}
\label{chap.arquitetura}

\section{Visão Geral}

A arquitetura final descreve um sistema de interação multimodal que permite aos utilizadores controlar a interface web do Google Maps através da combinação natural de comandos de voz e gestos. Utiliza  o protocolo \textbf{MMI} (MultiModal Interaction) para fazer a comunicação entre múltiplos componentes. O objetivo principal é permitir controlar a interface web (Google Maps) através da fusão de comandos de voz (Microfone) e gestos (Kinect), utilizando um modelo de integração centrado num motor de fusão.


\section{Componentes do Sistema}

A arquitetura pode ser dividida em quatro principais camadas que refletem o fluxo natural de processamento de informação: desde a captura de inputs do utilizador (perceção), passando pela interpretação e combinação desses inputs (fusão), seguida pela execução de ações (lógica de aplicação), e terminando com a apresentação de resultados ao utilizador (feedback). 

Nas subsecções seguintes, cada camada será descrita em detalhe, explicando os componentes em si bem como a função no sistema.

\subsection{Input Modalities}

A parte referente à modalidade gestual, ou seja, à captura de movimento é realizada pelo sensor Kinect. No fundo, os dados são processados pelo componente \textbf{GenericGesturesModality}, responsável por identificar os gestos do utilizador e enviá-los para o sistema de fusão.  

Por outro lado, a modalidade vocal, reconhecimento de voz, é processado pelo \textbf{Rasa NLU}. Este componente recebe o texto, após reconhecimento de fala, e extrai as intenções semânticas e entidades, enviando eventos estruturados para o sistema.


\subsection{Fusion \& Management}

O sistema de fusão é composto  por múltiplos componentes que comunicam através do protocolo MMI (MultiModal Interaction). O sistema inclui um \textbf{FusionEngine} baseado em SCXML que recebe eventos das duas modalidades: (1) comandos de voz, e (2) gestos. No fundo, ele encaminha esses eventos para o \textbf{Python Assistent} através da framework \textbf{MMI}.  


\subsection{Lógica de Aplicação}

O componente que atua como servidor dentral de lógica é o \textbf{Python Assistante}, que tem como função receber os comandos previamente fundidos pelo FusionEngine e traduzir os tais comandos em ações concretas. Após o nosso assistente fazer a tradução passa a responsabilidade para o Selenium web driver, que é este o responsável pela manipulação diretamente na interface web do Google Maps, executados as ações ordenadas pelo utilizador. 


\subsection{Feedback}

Como ultimo componente temos a camada de feedback onde está um serviço a rodar na porta 8083, o \textbf{servidor TTS}, que é responsável por gerar o feedback para o utilizador. Este servidor que comunica com uma \textbf{WebApp} baseada em browser, utilizando a Web Speech API para síntese de voz em português

\newpage

\section{Fluxo de Arquitetura}
\label{sec:fluxo_arquitetura}

Esta secção descreve o fluxo completo de processamento de dados no sistema, desde a captura de input multimodal até à execução de ações e geração de feedback ao utilizador. \\

\noindent O fluxo inicia-se com a captura simultânea de duas modalidades de input, sendo elas:

\begin{itemize}
    \item \textbf{Input Gestual}: Usa o sensor Kinect que captura os movimentos do utilizador. Os dados são enviados para o \textbf{Gesture Builder}, que processa e interpreta os gestos. Os gestos reconhecidos são então enviados para o \textbf{GenericGesturesModality}.

    \item \textbf{Input Vocal:} Usa o microfone que faz captura a voz do utilizador. Após o reconhecimento da fala, o texto resultante é processado pelo \textbf{Rasa NLU}, que realiza a análise de linguagem natural para extrair resultados.
\end{itemize}

Após esta fase de captura, os eventos provenientes das duas modalidades então convergem para o  \textbf{Fusion Engine}, o componente central da arquitetura, que é implementado em SCXML e é responsável por realizar a fusão dos comandos recolhidos e gerar comandos fundidos que são enviados para o \textbf{MMI}. \\

De seguida temos o componente \textbf{MMI} que no fundo atua como middleware de comunicação, recebendo os comandos previamente fundidos e envia-os para o \textbf{Assistant}. É neste mesmo \textbf{Assistant} onde os comando são mapeados para as ações concretas na interface web do Google Maps. Esta manipulação é feita usando o \textbf{Selenium}. \\

Como última fase da nossa arquitetura temos a fase de feedback multimodal ao utilizador. Para parte visual, o utilizador consegue ter feedback visual imediatamente através das alteraçoes na interface web do Google Maps. Já para o feedback sonoro, o \textbf{Assistant} envia mensagens de confirmação e informação para o \textbf{servidor TTS}. \\

Este fluxo contínuo permite uma interação natural e fluida, onde o utilizador pode combinar as interações de gestos e comandos de voz para controlar o Google Maps, recebendo feedback imediato das suas ações tanto visual como auditivamente. Na Figura \ref{fig:fusion-architecture} conseguimos ver de uma modo geral as interações entre os principais componentes do nosso sistema.

%%%%%%%%%%%%%%%%%%%%%%%%%%%%%%%%%%%%%%%%%
%     Adicionar imagem arquitetura      %
%%%%%%%%%%%%%%%%%%%%%%%%%%%%%%%%%%%%%%%%%

\begin{figure}[h]
    \centering
    \includegraphics[width=0.9\linewidth]{Images/Fusion-IM.png}
    \caption{Arquitetura de Fusão}
    \label{fig:fusion-architecture}
\end{figure}



% ============================================================================
% Chapter 3: Requisitos do Sistema
% ============================================================================


\chapter{Requisitos do Sistema}

O desenvolvimento e execução do sistema multimodal para controlo de mapas digitais requer um conjunto de componentes de hardware e software que garantam o correto funcionamento de todas as funcionalidades. Este capítulo apresenta os requisitos técnicos necessários para a instalação e operação do sistema.

\section{Requisitos de Hardware}

Para a execução adequada do sistema, é necessário dispor de equipamento capaz de suportar a captura de gestos e o processamento simultâneo de múltiplos componentes. Os requisitos de hardware incluem:

\begin{itemize}
    \item Sensor \textbf{Microsoft Kinect v2} para captura de movimentos e gestos. O sensor Kinect é essencial para a modalidade gestual do sistema, sendo responsável pela captura de movimentos corporais que são posteriormente processados e convertidos em comandos. ;
    \item Conexão estável à \textbf{Internet} para acesso ao Google Maps.
\end{itemize}

\section{Requisitos de Software}

O sistema foi desenvolvido para operar em ambiente Windows, aproveitando o suporte nativo do SDK do Kinect para esta plataforma. Os requisitos de software base incluem:

\begin{itemize}
    \item Sistema operativo \textbf{Windows 10} (64-bit);
    \item \textbf{Python} versão 3.10 ou superior;
    \item \textbf{Java JDK} versão 8 ou superior;
    \item \textbf{Google Chrome} ou Microsoft Edge (automação Selenium);
    \item \textbf{Kinect for Windows SDK v2};
\end{itemize}

\section{Bibliotecas e Frameworks Python}

O sistema utiliza diversas bibliotecas Python. As principais dependências incluem:

\begin{itemize}
    \item \textbf{RASA} versão 3.5.x para processamento de linguagem natural
    \item \textbf{Selenium} versão 4.15 ou superior para automação do browser
    \item \textbf{WebSockets} versão 12.0 ou superior para comunicação assíncrona
    \item \textbf{AsyncIO} para suporte de programação assíncrona (incluído no Python 3.10)
    \item \textbf{Colorlog} para sistema de logging estruturado
\end{itemize}

Todas as dependências Python estão especificadas no ficheiro \texttt{requirements.txt} incluído no projeto.

\section{Componentes do Framework Multimodal}

O sistema integra componentes da framework, responsáveis pela gestão de eventos e fusão de modalidades. Os componentes necessários são:

\begin{itemize}
    \item \textbf{MMI Framework} - implementação do protocolo W3C para comunicação multimodal
    \item \textbf{FusionEngine} - motor SCXML para fusão de eventos de diferentes modalidades
    \item \textbf{GenericGesturesModality} - módulo de processamento de gestos do Kinect
    \item \textbf{Keystores SSL} (keystore2.jks) para comunicação segura HTTPS/WSS
\end{itemize}

Estes componentes Java comunicam entre si através do protocolo MMI, garantindo sincronização temporal e fusão adequada de comandos vocais e gestuais. A configuração destes componentes é realizada através de ficheiros XML específicos incluídos no projeto.


%================================================================================

\begin{comment}

\chapter{Requisitos do Sistema}

O desenvolvimento de um sistema multimodal para controlo de mapas digitais requer a definição rigorosa de requisitos que garantam a implementação correta de todas as funcionalidades e a integração eficaz entre as diferentes modalidades. Este capítulo apresenta os requisitos funcionais, não-funcionais e técnicos que orientaram o desenvolvimento do sistema.

\section{Requisitos Funcionais}

Os requisitos funcionais definem as capacidades que o sistema deve proporcionar aos utilizadores através das diferentes modalidades de interação.

\subsection{Modalidade de Voz}

O sistema deve suportar reconhecimento e processamento de comandos de voz em português europeu (pt-PT), incluindo:

\begin{itemize}
    \item \textbf{Pesquisa e Navegação}
    \begin{itemize}
        \item Pesquisa de localizações, pontos de interesse e categorias (e.g., "procurar restaurantes em Lisboa")
        \item Obtenção de direções entre dois pontos com suporte para origem e destino
        \item Seleção de modo de transporte (carro, transportes públicos, a pé, bicicleta)
        \item Início e paragem de navegação turn-by-turn
        \item Seleção de rotas alternativas
    \end{itemize}

    \item \textbf{Controlo do Mapa}
    \begin{itemize}
        \item Operações de zoom (aumentar/diminuir) com níveis variáveis (muito, pouco)
        \item Alteração do tipo de vista (satélite, terreno, mapa padrão)
        \item Recentragem do mapa na localização atual do utilizador
        \item Centragem em localizações específicas
        \item Ativação e desativação da camada de tráfego
    \end{itemize}

    \item \textbf{Informação de Locais}
    \begin{itemize}
        \item Visualização de detalhes de estabelecimentos
        \item Consulta de avaliações e fotografias
        \item Obtenção de horários de funcionamento
        \item Verificação de informações gerais sobre localizações
        \item Descoberta de pontos de interesse numa área ("o que há aqui?")
    \end{itemize}

    \item \textbf{Informação de Viagem}
    \begin{itemize}
        \item Consulta de duração estimada da viagem
        \item Consulta de distância total do percurso
        \item Mudança de modo de transporte durante planeamento
        \item Alternância entre rotas disponíveis
    \end{itemize}

    \item \textbf{Seleção e Confirmação}
    \begin{itemize}
        \item Seleção de resultados de pesquisa por ordem ou nome
        \item Confirmação de ações (afirmação/negação)
        \item Cancelamento de operações em curso
    \end{itemize}

    \item \textbf{Interação Geral}
    \begin{itemize}
        \item Comandos de saudação e despedida
        \item Pedidos de ajuda sobre funcionalidades disponíveis
        \item Agradecimentos e feedback ao sistema
    \end{itemize}
\end{itemize}

\subsection{Modalidade de Gestos}

O sistema deve reconhecer e processar gestos capturados pelo sensor Kinect v2, nomeadamente:

\begin{itemize}
    \item \textbf{Gestos de Navegação}
    \begin{itemize}
        \item Swipe (deslizar) em quatro direções: cima, baixo, esquerda, direita
        \item Zoom in/out através de aproximação/afastamento das mãos
    \end{itemize}

    \item \textbf{Gestos de Seleção}
    \begin{itemize}
        \item Polegar para cima (thumbs up) para afirmação
        \item Braço direito para cima para ver opção de cima
        \item Braço esquerdo para baixo para ver opção de baixo
    \end{itemize}

    \item \textbf{Gestos de Filtros e Funcionalidades}
    \begin{itemize}
        \item Ativação de filtro de hotéis
        \item Ativação de filtro de restaurantes
        \item Ativação de filtro de transportes
        \item Acesso a "Coisas a fazer" (pontos turísticos)
    \end{itemize}

    \item \textbf{Gestos de Street View}
    \begin{itemize}
        \item Entrada em modo Street View
        \item Saída de modo Street View
    \end{itemize}

    \item \textbf{Gestos de Navegação em Listas}
    \begin{itemize}
        \item Subir/descer em listas de resultados
        \item Seleção de opções apresentadas
    \end{itemize}
\end{itemize}

\subsection{Fusão Multimodal}

O motor de fusão deve proporcionar as seguintes capacidades:

\begin{itemize}
    \item \textbf{Combinação Temporal}: Fusão de eventos de diferentes modalidades que ocorrem num intervalo temporal configurável (janela de fusão)
    \item \textbf{Resolução de Ambiguidades}: Desambiguação de comandos através do contexto fornecido por múltiplas modalidades
    \item \textbf{Priorização Contextual}: Seleção da modalidade mais apropriada baseada no tipo de comando e contexto de utilização
    \item \textbf{Validação de Comandos}: Verificação de consistência entre modalidades antes da execução
    \item \textbf{Tratamento de Conflitos}: Resolução de comandos contraditórios através de regras de priorização
\end{itemize}

\subsection{Controlo de Google Maps}

O módulo de execução deve ser capaz de:

\begin{itemize}
    \item Automatizar as principais operações de interface do Google Maps através de Selenium WebDriver
    \item Executar pesquisas e processar resultados com múltiplas entradas
    \item Manipular controlos de navegação
    \item Aceder e extrair informação de painéis laterais (detalhes, avaliações, fotografias)
    \item Gerir sessões persistentes do browser sem perda de contexto
    \item Recuperar de erros de carregamento e elementos não encontrados
\end{itemize}

\subsection{Feedback Text-to-Speech}

O sistema de feedback áudio deve:

\begin{itemize}
    \item Sintetizar fala em português europeu (pt-PT) com voz natural
    \item Fornecer confirmação de ações executadas com sucesso
    \item Reportar erros e problemas de forma clara e compreensível
    \item Informar sobre resultados de pesquisas e consultas
    \item Solicitar confirmação quando necessário
    \item Fornecer orientação contextual ao utilizador
\end{itemize}

\section{Requisitos Não-Funcionais}

Os requisitos não-funcionais estabelecem critérios de qualidade e desempenho do sistema.

\subsection{Desempenho e Latência}

\begin{itemize}
    \item Latência de reconhecimento de voz inferior a 2 segundos (desde captura até processamento NLU)
    \item Latência de reconhecimento de gestos inferior a 500ms (desde captura até evento MMI)
    \item Tempo de processamento de fusão inferior a 200ms
    \item Execução de comandos no mapa em menos de 1 segundo (exceto operações de rede)
    \item Tempo de resposta TTS inferior a 500ms após geração de mensagem
\end{itemize}

\subsection{Precisão e Confiabilidade}

\begin{itemize}
    \item Taxa de reconhecimento de intents (RASA) superior a 85\% para comandos bem formulados
    \item Taxa de reconhecimento de gestos superior a 80\% em condições de iluminação adequadas
    \item Sistema de confiança (confidence score) para comandos ambíguos:
    \begin{itemize}
        \item Alta confiança ($\geq$ 0.80): execução imediata
        \item Média confiança (0.60--0.79): solicitação de confirmação
        \item Baixa confiança (< 0.60): rejeição com pedido de repetição
    \end{itemize}
    \item Recuperação automática de falhas de conexão WebSocket
    \item Retry logic para operações Selenium com até 3 tentativas
\end{itemize}

\subsection{Usabilidade}

\begin{itemize}
    \item Interface multimodal intuitiva sem necessidade de treino prévio
    \item Comandos de voz em linguagem natural portuguesa coloquial
    \item Gestos baseados em movimentos naturais e culturalmente reconhecíveis
    \item Feedback imediato e claro de todas as ações e estados
    \item Mensagens de erro descritivas e orientações para correção
    \item Sistema de ajuda contextual acessível por voz
\end{itemize}

\subsection{Robustez e Tolerância a Falhas}

\begin{itemize}
    \item Tratamento robusto de erros de reconhecimento de voz e gestos
    \item Validação de entidades extraídas antes da execução
    \item Gestão de timeouts em todas as operações de rede
    \item Logging detalhado e estruturado para debugging
    \item Recuperação graceful de crashes do browser
    \item Reinício automático de componentes em caso de falha crítica
\end{itemize}

\subsection{Manutenibilidade e Extensibilidade}

\begin{itemize}
    \item Arquitetura modular baseada em componentes independentes
    \item Separação clara entre lógica de negócio e integração
    \item Configuração centralizada através de ficheiros de settings
    \item Documentação inline de código e APIs
    \item Estrutura de handlers extensível para novos intents
    \item Sistema de plugins para adicionar novas modalidades
\end{itemize}

\subsection{Portabilidade}

\begin{itemize}
    \item Instalação simplificada através de ambiente virtual Python
    \item Dependências claramente especificadas em requirements.txt
    \item Scripts de inicialização automatizados (start.bat)
    \item Caminhos relativos para portabilidade entre máquinas
    \item Documentação completa de setup para novos utilizadores
\end{itemize}

\section{Requisitos Técnicos}

Os requisitos técnicos especificam as tecnologias, protocolos e infraestrutura necessários.

\subsection{Hardware}

\begin{itemize}
    \item \textbf{Sensor Kinect}: Microsoft Kinect v2 com SDK apropriado
    \item \textbf{Conectividade}: Porta USB 3.0 dedicada para o Kinect
    \item \textbf{Microfone}: Sistema de captura de áudio (integrado no Kinect ou externo)
    \item \textbf{Processador}: CPU multi-core (recomendado Intel i5 ou superior)
    \item \textbf{Memória}: Mínimo 8GB RAM para execução simultânea de todos os componentes
    \item \textbf{Rede}: Conexão à Internet para acesso ao Google Maps e APIs
\end{itemize}

\subsection{Software e Plataformas}

\begin{itemize}
    \item \textbf{Sistema Operativo}: Windows 10/11 (64-bit)
    \item \textbf{Python}: Versão 3.10 (compatibilidade com RASA e TensorFlow)
    \item \textbf{Java}: JDK 8 ou superior para MMI Framework e FusionEngine
    \item \textbf{Browser}: Google Chrome ou Microsoft Edge (automação Selenium)
    \item \textbf{ChromeDriver}: Versão compatível com browser instalado (gestão automática via Selenium Manager)
\end{itemize}

\subsection{Bibliotecas e Frameworks Python}

\begin{itemize}
    \item \textbf{RASA}: Framework de NLU versão 3.5.x para processamento de linguagem natural
    \item \textbf{Selenium}: Biblioteca de automação web versão 4.15+ para controlo do browser
    \item \textbf{WebSockets}: Biblioteca websockets 12.0+ para comunicação assíncrona
    \item \textbf{AsyncIO}: Suporte para programação assíncrona (incluído no Python 3.10)
    \item \textbf{Pytest}: Framework de testes automatizados (desenvolvimento)
    \item \textbf{Colorlog}: Logging colorido para melhor debugging
\end{itemize}

\subsection{Componentes Java}

\begin{itemize}
    \item \textbf{MMI Framework}: Implementação do protocolo W3C MMI para comunicação multimodal
    \item \textbf{FusionEngine}: Motor SCXML para orquestração e fusão de modalidades
    \item \textbf{GenericGesturesModality}: Módulo de processamento de gestos Kinect
    \item \textbf{Keystores SSL}: Certificados keystore2.jks para comunicação segura HTTPS/WSS
\end{itemize}

\subsection{Protocolos e Padrões}

\begin{itemize}
    \item \textbf{MMI Protocol (W3C)}: Protocolo padrão para comunicação entre modalidades e motor de fusão
    \item \textbf{WebSocket}: Protocolo bidirecional para comunicação em tempo real
    \item \textbf{HTTPS/WSS}: Protocolos seguros para todas as comunicações de rede
    \item \textbf{JSON}: Formato de serialização de mensagens e dados
    \item \textbf{XML}: Formato de mensagens MMI e configuração SCXML
    \item \textbf{SCXML}: State Chart XML para definição de máquinas de estado de fusão
\end{itemize}

\subsection{Arquitetura de Comunicação}

O sistema utiliza múltiplas portas de comunicação para diferentes componentes:

\begin{itemize}
    \item \textbf{Porta 8000}: MMI Framework Client (comunicação com FusionEngine)
    \item \textbf{Porta 8001}: FusionEngine HTTPS (receção de eventos de modalidades)
    \item \textbf{Porta 8005}: MMI Server WSS (distribuição de comandos fundidos)
    \item \textbf{Porta 8082}: WebApp HTTPS Server (interface web e captura de voz)
    \item \textbf{Porta 8083}: TTS WebSocket Server (distribuição de mensagens TTS)
    \item \textbf{Porta 5005}: RASA API Server (processamento NLU)
\end{itemize}

\subsection{Estrutura de Dados}

\subsubsection{Formato de Eventos MMI}

Os eventos seguem o protocolo MMI com estrutura XML encapsulando JSON:

\begin{verbatim}
<mmi:mmi xmlns:mmi="http://www.w3.org/2008/04/mmi-arch">
  <mmi:startRequest mmi:context="..." mmi:source="..."
                    mmi:target="..." mmi:requestID="...">
    <mmi:data>
      <command>
        {
          "recognized": ["SPEECH"],
          "nlu": {
            "intent": {"name": "search_location", "confidence": 0.92},
            "entities": [{"entity": "location", "value": "Lisboa"}],
            "text": "procurar Lisboa"
          }
        }
      </command>
    </mmi:data>
  </mmi:startRequest>
</mmi:mmi>
\end{verbatim}

\subsubsection{Estrutura de Resposta RASA}

\begin{verbatim}
{
  "intent": {
    "name": "get_directions",
    "confidence": 0.87
  },
  "entities": [
    {"entity": "destination", "value": "Porto", "start": 8, "end": 13},
    {"entity": "transport_mode", "value": "carro", "start": 17, "end": 22}
  ],
  "text": "ir para Porto de carro"
}
\end{verbatim}
\end{comment}



% ============================================================================
% Chapter 4: Arquitetura de Fusão Multimodal
% ============================================================================

\chapter{Arquitetura de Fusão Multimodal}


Este capítulo detalha a implementação do motor de fusão multimodal, componente central que permite ao sistema combinar informação proveniente das modalidades de voz e gestos. O motor de fusão implementa diferentes estratégias de integração modal, permitindo que o utilizador interaja com o sistema de forma natural e flexível, utilizando uma ou ambas as modalidades conforme a tarefa e contexto.

\section{Conceitos Fundamentais de Fusão}

A fusão multimodal é o processo de combinar informação proveniente de múltiplas modalidades de input para gerar comandos unificados que são enviados para a aplicação, no âmbito deste projeto era a fusão entre as modalidades de interação de gestos e voz:

\begin{itemize}
    \item \textbf{Modalidade de Voz}: Comandos vindos do RASA NLU, contendo intents e entidades extraídas de linguagem natural;
    \item \textbf{Modalidade de Gestos}: Comandos recolhidos pelo Kinect v2 através do GenericGesturesModality, representando gestos físicos do utilizador.
\end{itemize}


\section{Implementação de Complementaridade}

A fusão complementar ocorre quando as duas modalidades fornecem informação distinta que, quando combinada. 

O sistema implementa 8 combinações complementares que permitem ao utilizador controlar simultaneamente zoom e direção de navegação. Estas combinações cobrem todas as possibilidades de zoom (in/out) com as quatro direções de movimento (up/down/left/right). A Tabela \ref{tab:complementary-fusion} indicas as possiveis formas de combinação de comandos de voz com os gestos gerando apenas um comando de saída. 

\subsection{Combinações Implementadas}

\begin{longtable}{p{0.30\textwidth} p{0.30\textwidth} p{0.33\textwidth}}
    \caption{Comandos de fusão complementar implementados} 
    \label{tab:complementary-fusion} \\
    
    % --- Cabeçalho da Primeira Página ---
    \toprule
    \textbf{Comando de Voz} & \textbf{Gesto} & \textbf{Comando Fundido} \\
    \midrule
    \endfirsthead
    \endlastfoot
    
    % --- Redundância (3 comandos) ---
    \multicolumn{3}{l}{\textbf{Redundância}} \\
    \midrule
    \texttt{Aproximar} & Zoom In (expandir braços) & ZOOM\_IN  \\
    \texttt{Afastar} & Zoom Out (expandir braços) & ZOOM\_OUT \\
    \texttt{Cancelar} & Exit Street & CANCEL \\
    \midrule
    
    % --- Fusão Complementar Original (8 comandos) ---
    \multicolumn{3}{l}{\textbf{Fusão Complementar - Zoom + Direções}} \\
    \midrule
    \texttt{Aproximar (ZOOM\_IN)} & Swipe Up & ZOOM\_IN\_UP \\ 
    \texttt{Aproximar (ZOOM\_IN)} & Swipe Down & ZOOM\_IN\_DOWN \\ 
    \texttt{Aproximar (ZOOM\_IN)} & Swipe Left & ZOOM\_IN\_LEFT \\ 
    \texttt{Aproximar (ZOOM\_IN)} & Swipe Right & ZOOM\_IN\_RIGHT \\ 
    \texttt{Afastar (ZOOM\_OUT)} & Swipe Up & ZOOM\_OUT\_UP \\ 
    \texttt{Afastar (ZOOM\_OUT)} & Swipe Down & ZOOM\_OUT\_DOWN \\ 
    \texttt{Afastar (ZOOM\_OUT)} & Swipe Left & ZOOM\_OUT\_LEFT \\ 
    \texttt{Afastar (ZOOM\_OUT)} & Swipe Right & ZOOM\_OUT\_RIGHT \\ 
    \texttt{Obter direções (GET\_DIRECTIONS)} & Transports (gesto) & DIRECTIONS\_PUBLIC\_TRANSPORT \\
    \texttt{Mostrar tráfego (SHOW\_TRAFFIC)} & Zoom Out (gesto) & TRAFFIC\_OVERVIEW \\
    Restaurants (gesto) & \texttt{Centrar mapa (RECENTER\_MAP)} & FILTER\_RESTAURANTS\_CENTER \\
    Forward (gesto) & \texttt{Iniciar navegação (START\_NAVIGATION)} & STREET\_VIEW\_FORWARD\_CONTINUOUS \\
    \bottomrule
\end{longtable}

A fusão redundante ocorre quando ambas as modalidades expressam a mesma intenção, per-
mitindo ao utilizador escolher a forma de interação mais conveniente para cada situação. Este
paradigma aumenta a robustez do sistema.

A fusão redundante permite que o utilizador escolha livremente entre modalidades e possa executar um comando com esssa modalidade escolhida. No fundo, o \textbf{FusionEngine} reconhece a redundância e gere um único evento fundido ambas as modalidades, e desta forma, apesar de para a mesma ação haver possibilidade de ambas as modalidades, o \textbf{Assistant} executa a ação apenas uma vez. O sistema implementa redundância para três comandos críticos.

\subsection{Comandos Redundantes Implementados}


\begin{longtable}{p{0.30\textwidth} p{0.30\textwidth} p{0.33\textwidth}}
    \caption{Comandos de fusão redundante implementados} 
    \label{tab:redundant-fusion} \\
    
    % --- Cabeçalho da Primeira Página ---
    \toprule
    \textbf{Comando de Voz} & \textbf{Gesto} & \textbf{Resultado} \\
    \midrule
    \endfirsthead
    \endlastfoot
    
    \texttt{Aproximar} & Zoom In (expandir braços)& ZOOM\_IN  \\
    \texttt{Afastar} & Zoom Out (expandir braços) & ZOOM\_OUT \\
    \texttt{Cancelar} & Exit Street & CANCEL \\
\end{longtable}



\section{Fluxo Completo de Fusão}

O processamento completo de um comando multimodal segue o seguinte fluxo:

\begin{enumerate}
    \item \textbf{Captura}: Modalidades capturam input do utilizador (voz via microfone, gestos via Kinect)

    \item \textbf{Reconhecimento}:
    \begin{itemize}
        \item Voz processada por Speech-to-Text + RASA NLU
        \item Gestos processados por GenericGesturesModality
    \end{itemize}

    \item \textbf{Envio MMI}: Cada modalidade envia evento MMI para FusionEngine (porta 8001):
    \begin{itemize}
        \item Voz: \texttt{POST /IM/USER1/SPEECHIN}
        \item Gestos: \texttt{POST /IM/USER1/GESTURES}
    \end{itemize}

    \item \textbf{Fusão SCXML}: FusionEngine processa eventos:
    \begin{itemize}
        \item Armazena primeiro evento e inicia janela temporal
        \item Se segundo evento chegar, funde modalidades
        \item Se timeout, processa evento isoladamente
    \end{itemize}

    \item \textbf{Envio para Assistant}: FusionEngine envia evento fundido via MMI Server (porta 8005)

    \item \textbf{Extração NLU}: Assistant extrai intent através de \texttt{nlu\_extractor()}:
\begin{verbatim}
if command_type == "FUSION":
    fusion_intent = recognized[1]  # e.g., "ZOOM_IN_UP"
    intent = fusion_intent
\end{verbatim}

    \item \textbf{Routing}: IntentRouter encaminha para handler apropriado

    \item \textbf{Execução}: Handler executa ações no Google Maps via Selenium

    \item \textbf{Feedback}: Sistema fornece feedback auditivo via TTS
\end{enumerate}

\subsection{Exemplo Completo: Fusão Complementar}

\noindent \textbf{Cenário:} Utilizador diz "aproximar" e faz swipe up simultaneamente.

\begin{enumerate}
    \item Web Speech API captura "aproximar" $\rightarrow$ RASA extrai intent \texttt{zoom\_in}
    \item WebApp envia \texttt{[SPEECH][ZOOM\_IN]} para FusionEngine
    \item FusionEngine entra em estado \texttt{sspeech[zoom\_in]}
    \item Kinect deteta swipe up $\rightarrow$ GenericGesturesModality envia \texttt{[GESTURES][SWIPE\_UP]}
    \item FusionEngine deteta transição complementar, cancela timer
    \item Estado de fusão combina \texttt{data1} (voz) + \texttt{data2} (gesto)
    \item FusionEngine envia \texttt{[FUSION][ZOOM\_IN\_UP]} para Assistant
    \item Assistant identifica comando fundido, invoca \texttt{ZoomInUpHandler}
    \item Handler executa: \texttt{zoom\_in(2)} $\rightarrow$ pausa 300ms $\rightarrow$ \texttt{pan\_map("up", 3)}
    \item TTS anuncia: "A aproximar e mover para cima"
\end{enumerate}

\section{Melhorias e Otimizações Implementadas}

Durante a fusão e integração dos dois projetos (sistema multimodal de reconhecimento de voz/gestos e controlo automatizado do Google Maps), foram implementadas diversas melhorias significativas a nível de lógica, treino de modelos, tratamento de erros, desempenho e experiência de utilizador. Esta secção detalha as principais otimizações realizadas.

\subsection{Melhorias no Processamento de Linguagem Natural}

\subsubsection{Expansão do Corpus de Treino RASA}

O modelo RASA foi significativamente melhorado através da expansão do corpus de treino:

\begin{itemize}
    \item \textbf{Aumento de Exemplos}: Adição de mais de 100 novos exemplos de treino para intents críticos, com foco em variações linguísticas do português europeu coloquial
    
    \item \textbf{Modos de Transporte}: Expansão massiva do intent \texttt{change\_transport\_mode} para suportar formas isoladas:
    \begin{itemize}
        \item Palavras isoladas: "pé", "bicicleta", "carro", "autocarro", "comboio", "metro"
        \item Formas verbais: "caminhar", "andar", "conduzir"
        \item Variações compostas: "ir a pé", "ir de bicicleta", "ir de carro"
        \item Comandos completos: "mudar para transportes públicos", "alterar modo de transporte"
    \end{itemize}
    
\end{itemize}

\subsubsection{Mapeamento Semântico de Entidades}

Criação de sistema robusto de mapeamento para entidades de transporte em \texttt{enums.py}:

\begin{itemize}
    \item \textbf{Mapeamento Bilingue}: Suporte para termos em inglês e português:
    \begin{itemize}
        \item \texttt{TransportMode.DRIVING}: "driving", "drive", "car", "carro", "conduzir", "automóvel", "viatura"
        \item \texttt{TransportMode.WALKING}: "walking", "walk", "on foot", "foot", "pé", "a pé", "andar", "caminhar", "caminhando"
        \item \texttt{TransportMode.TRANSIT}: "transit", "public transport", "bus", "train", "metro", "subway", "transportes públicos", "transportes", "público", "públicos", "autocarro", "comboio", "trem"
        \item \texttt{TransportMode.CYCLING}: "cycling", "bike", "bicycle", "biking", "bicicleta", "bici", "ciclismo"
    \end{itemize}
    
    \item \textbf{Normalização Automática}: Conversão case-insensitive e remoção de espaços redundantes antes de lookup
    
    \item \textbf{Fallback Hierárquico}: Sistema de três níveis de reconhecimento:
    \begin{enumerate}
        \item Matching direto no dicionário de mapeamento
        \item Comparação com valores enum exatos
        \item Lançamento de \texttt{ValueError} descritivo se não reconhecido
    \end{enumerate}
\end{itemize}

\subsection{Melhorias no Tratamento de Erros}

\subsubsection{Inferência de Texto para Entidades Não Extraídas}

Implementação de sistema de fallback inteligente em \texttt{trip\_info\_handler.py} quando RASA reconhece o intent mas não extrai a entidade:

\begin{verbatim}
# Se RASA não extraiu entidade, inferir do texto completo
if not transport_mode_str and context.metadata:
    text = context.metadata.get("text", "").strip().lower()
    mode_mappings = {
        "caminhar": "pé",
        "andar": "pé",
        "pé": "pé",
        "bicicleta": "bicicleta",
        "carro": "carro",
        "transportes": "transportes públicos",
        # ... mais mapeamentos
    }
    
    for keyword, mode in mode_mappings.items():
        if keyword in text:
            transport_mode_str = mode
            logger.info(f"Inferred '{mode}' from text '{text}'")
            break
\end{verbatim}

\textbf{Impacto}: Eliminou loops infinitos onde o sistema pedia repetidamente "Que meio de transporte?" mesmo quando o utilizador já tinha especificado (e.g., "caminhar").

\subsubsection{Mensagens de Erro Amigáveis}

Substituição de mensagens técnicas por feedback em português claro:

\begin{itemize}
    \item \textbf{Antes}: \texttt{"I don't know how to handle 'nlu\_fallback'"}
    \item \textbf{Depois}: \texttt{"Desculpa, não entendi o que pediste."}
\end{itemize}

Implementação em \texttt{intent\_router.py}:

\begin{verbatim}
if handler is None:
    return IntentResponse(
        success=False,
        message="Desculpa, não entendi o que pediste.",
        data={"error": f"No handler for intent '{context.intent}'"}
    )
\end{verbatim}

\subsubsection{Remoção de Restrições Desnecessárias}

Eliminação de verificação restritiva que impedia mudança de modo de transporte durante planeamento de rota:

\begin{itemize}
    \item \textbf{Problema Anterior}: Sistema exigia painel de direções completamente carregado antes de permitir mudança de transporte
    \item \textbf{Solução}: Remoção do check, permitindo mudança durante qualquer fase do planeamento
    \item \textbf{Resultado}: Utilizador pode mudar modo de transporte assim que chega à página de planeamento de rotas
\end{itemize}

\subsection{Otimizações de Desempenho}

\subsubsection{Redução de Latências em Operações de Pesquisa}

Otimização agressiva de timeouts e sleeps em \texttt{search\_handler.py}:

\begin{table}[h]
\centering
\caption{Redução de latências em operações críticas}
\label{tab:performance-improvements}
\begin{tabular}{lccr}
\toprule
\textbf{Operação} & \textbf{Antes} & \textbf{Depois} & \textbf{Melhoria} \\
\midrule
Sleep após pesquisa & 1.0s & 0.3s & 70\% \\
Timeout de validação & 2.0s & 1.5s & 25\% \\
Timeout de resultados & 8.0s & 4.0s & 50\% \\
Timeout de painel & 5.0s & 3.0s & 40\% \\
\midrule
\textbf{Total Estimado} & \textbf{16.0s} & \textbf{8.8s} & \textbf{45\%} \\
\bottomrule
\end{tabular}
\end{table}

\textbf{Impacto Real}: Comando "ir para Aveiro" reduziu de 10--15 segundos para 5--8 segundos (melhoria de aproximadamente 50\%).

\subsubsection{Otimização de Sleeps em Handlers de Fusão}

Redução de delays desnecessários em \texttt{fusion\_handler.py}:

\begin{verbatim}
# Antes
time.sleep(0.3)  # Cada operação

# Depois  
time.sleep(0.1)  # Redução de 70%
\end{verbatim}

\textbf{Resultado}: Comandos de fusão complementar (e.g., zoom + swipe) executam 0.6s mais rápido.


\subsection{Melhorias na Experiência de Utilizador}

\subsubsection{Wake Word Detection com Web Speech API}

Implementação de sistema de ativação por palavra-chave "maps":

\begin{itemize}
    \item \textbf{Reconhecimento Contínuo}: \texttt{wakeWordRecognition} escuta continuamente em background
    \item \textbf{Ativação por Voz}: Deteta "maps", "mapa" ou "map" e ativa reconhecimento de comandos
    \item \textbf{Feedback Visual}: Círculo verde indica sistema pronto para receber comando
    \item \textbf{Gestão de Estado}: Pausa wake word durante TTS para evitar auto-ativação
\end{itemize}

\subsubsection{Auto-Ativação Após Perguntas do Sistema}

Solução para problema crítico de UX onde utilizador tinha de repetir wake word para responder a perguntas:

\begin{itemize}
    \item \textbf{Problema Anterior}: Sistema pergunta "Devo iniciar a navegação?" mas utilizador tem de dizer "maps" antes de responder "sim"
    
    \item \textbf{Solução Implementada}:
    \begin{enumerate}
        \item Deteção automática de mensagens TTS terminadas com "?"
        \item Flag \texttt{pendingQuestionActivation} controla estado
        \item Handler \texttt{handleTTSEnd()} ativa reconhecimento automaticamente após TTS terminar
        \item Sistema de fallback baseado em duração estimada para browsers sem suporte a \texttt{msg.onend}
    \end{enumerate}
\end{itemize}

Implementação com proteção contra duplicação:

\begin{verbatim}
var ttsEndHandled = false;

function handleTTSEnd() {
    if (ttsEndHandled) return;  // Previne execução dupla
    ttsEndHandled = true;

    if (pendingQuestionActivation) {
        pendingQuestionActivation = false;
        setTimeout(function() {
            recognition.start();
            circle.animate(20, 0, 'now').attr({ fill: '#00a431' });
        }, 500);
    }
}

// Handler primário
msg.onend = handleTTSEnd;

// Fallback para browsers sem onend confiável
var estimatedDuration = (msg.text.length / 15) * 1000;
setTimeout(function() {
    if (!window.speechSynthesis.speaking && !ttsEndHandled) {
        handleTTSEnd();
    }
}, estimatedDuration + 1000);
\end{verbatim}

\subsection{Melhorias na Arquitetura e Manutenibilidade}

\subsubsection{Modularização de Intent Handlers}

Criação de handlers especializados para fusão complementar em \texttt{fusion\_handler.py}:

\begin{itemize}
    \item \texttt{ZoomInUpHandler}, \texttt{ZoomInDownHandler}, \texttt{ZoomInLeftHandler}, \texttt{ZoomInRightHandler}
    \item \texttt{ZoomOutUpHandler}, \texttt{ZoomOutDownHandler}, \texttt{ZoomOutLeftHandler}, \texttt{ZoomOutRightHandler}
    \item \texttt{DirectionsPublicTransportHandler}
    \item \texttt{TrafficOverviewHandler}
    \item \texttt{FilterRestaurantsCenterHandler}
    \item \texttt{StreetViewForwardContinuousHandler}
\end{itemize}

Cada handler encapsula lógica específica do comando fundido, facilitando manutenção e extensão.

\subsubsection{Logging Estruturado e Detalhado}

Implementação de logging colorido e estruturado usando \texttt{colorlog}:

\begin{verbatim}
logger.info(f"Inferred transport mode '{mode}' from text '{text}'")
logger.warning(f"Failed to click transport button: {str(e)}")
logger.error(f"No handler registered for intent '{intent}'")
\end{verbatim}

\textbf{Benefícios}:
\begin{itemize}
    \item Debugging rápido de problemas em produção
    \item Rastreabilidade de decisões de fusão e inferência
    \item Identificação de padrões de falha
\end{itemize}

\subsubsection{Gestão de Dependências}

Resolução de conflitos de versões Python/RASA:

\begin{itemize}
    \item \textbf{Ambiente Virtual Python 3.10}: Compatibilidade com RASA 3.5 e TensorFlow
    \item \textbf{Scripts de Inicialização Automática}: \texttt{start.bat} ativa ambiente correto
    \item \textbf{Requirements Modulares}: Separação entre dependências RASA (global) e projeto (venv)
\end{itemize}

\subsection{Expansão de Funcionalidades Multimodais}


\subsubsection{Geração SCXML Atualizada}

Atualização de \texttt{GenFusionSCXML.java} para gerar máquina de estados com novos comandos:

\begin{verbatim}
// New Complementary Fusion Commands (2026)
fg.Complementary(Speech.GET_DIRECTIONS, Gesture.TRANSPORTS, 
                 Output.DIRECTIONS_PUBLIC_TRANSPORT);
fg.Complementary(Speech.SHOW_TRAFFIC, Gesture.ZOOM_OUT, 
                 Output.TRAFFIC_OVERVIEW);
fg.Complementary(Gesture.RESTAURANTS, Speech.RECENTER_MAP, 
                 Output.FILTER_RESTAURANTS_CENTER);
fg.Complementary(Gesture.FORWARD, Speech.START_NAVIGATION, 
                 Output.STREET_VIEW_FORWARD_CONTINUOUS);
\end{verbatim}

\subsection{Impacto Global das Melhorias}

\subsubsection{Métricas de Desempenho}

\begin{table}[h]
\centering
\caption{Impacto quantitativo das otimizações}
\label{tab:global-impact}
\begin{tabular}{lcc}
\toprule
\textbf{Métrica} & \textbf{Antes} & \textbf{Depois} \\
\midrule
Latência média de pesquisa & 10--15s & 5--8s \\
Taxa de reconhecimento de transporte & 70\% & 95\% \\
Comandos de fusão disponíveis & 0 & 15 \\
Loops infinitos de confirmação & Frequente & Eliminado \\
Tempo de resposta a perguntas & 2--3 comandos & 1 comando \\
Falhas de locator Selenium & 15\% & $<$5\% \\
\bottomrule
\end{tabular}
\end{table}

\subsubsection{Melhorias Qualitativas}

\begin{itemize}
    \item \textbf{Naturalidade}: Wake word elimina necessidade de clicar botão antes de cada comando
    \item \textbf{Fluidez}: Auto-ativação após perguntas torna diálogo mais natural
    \item \textbf{Robustez}: Sistema tolera variações linguísticas e erros de extração de entidades
    \item \textbf{Inteligência}: Inferência de contexto reduz necessidade de comandos explícitos
    \item \textbf{Velocidade}: Reduções de 45--50\% em operações críticas melhoram perceção de responsividade
\end{itemize}

\subsection{Conclusão}

As melhorias implementadas transformaram o sistema de um protótipo funcional para uma aplicação robusta e utilizável em contextos reais. A combinação de otimizações de desempenho (redução de 50\% em latências), melhorias de UX (wake word, auto-ativação), robustez de NLU (inferência de texto, corpus expandido), e expansão de funcionalidades (4 novos comandos de fusão) resultou num sistema significativamente mais capaz e agradável de utilizar.

O trabalho realizado demonstra a importância de iterar sobre sistemas multimodais através de testes com utilizadores reais, identificando pontos de fricção e implementando soluções elegantes que respeitam os princípios de interação natural.



\begin{comment}
    
A fusão multimodal constitui o núcleo do sistema, sendo responsável por combinar informação proveniente de diferentes modalidades de input (voz e gestos) para criar uma experiência de interação coesa e natural. Este capítulo detalha a arquitetura e mecanismos de fusão implementados.

\section{Abordagem de Fusão}

O sistema implementa uma arquitetura de fusão híbrida baseada em SCXML (State Chart XML), combinando fusão baseada em regras com fusão temporal. Esta abordagem permite processar eficazmente tanto comandos unimodais (uma única modalidade) como multimodais (combinação de modalidades).

\subsection{Paradigma de Fusão}

O FusionEngine adota os seguintes princípios:

\begin{description}
    \item[Fusão Frame-Based] Utilização de frames semânticos para representar comandos, onde slots podem ser preenchidos por diferentes modalidades
    \item[Fusão Temporal] Combinação de eventos que ocorrem dentro de uma janela temporal configurável (por defeito 2 segundos)
    \item[Fusão por Complementaridade] Modalidades fornecem informação complementar (e.g., voz especifica ação, gesto especifica localização)
    \item[Fusão por Redundância] Múltiplas modalidades expressam a mesma informação, aumentando robustez
    \item[Fusão Hierárquica] Decisões de fusão organizadas em níveis (sintático, semântico, pragmático)
\end{description}

\subsection{Janela Temporal de Fusão}

O sistema mantém uma janela temporal deslizante de 2 segundos onde eventos de diferentes modalidades podem ser fundidos. O funcionamento é:

\begin{enumerate}
    \item Evento de modalidade A chega ao FusionEngine (t = 0ms)
    \item FusionEngine entra em estado COLLECTING e inicia timer de 2000ms
    \item Eventos adicionais que chegam antes do timeout são adicionados ao buffer
    \item Ao expirar o timeout ou receber evento especial de conclusão:
    \begin{itemize}
        \item Transita para estado FUSING
        \item Processa todos os eventos no buffer
        \item Gera comando fundido
    \end{itemize}
    \item Comando fundido é validado e enviado para execução
\end{enumerate}

\subsubsection{Exemplo de Fusão Temporal}

\begin{enumerate}
    \item \textbf{t=0ms}: Utilizador diz "aumenta o zoom aqui" \\
    RASA extrai: intent=\texttt{zoom\_in}, confidence=0.88

    \item \textbf{t=300ms}: Utilizador dá swipe para localização no mapa \\
    Kinect extrai: gesture=\texttt{select}, coordinates=(0.65, 0.42)

    \item \textbf{t=2000ms}: Timeout da janela temporal expira

    \item \textbf{Fusão}: FusionEngine combina os dois eventos: \\
    \texttt{\{action: "zoom\_in", target: \{x: 0.65, y: 0.42\}, source: "FUSION"\}}

    \item \textbf{Execução}: Assistant executa zoom centrado nas coordenadas especificadas
\end{enumerate}

\subsection{Estratégias de Resolução de Conflitos}

Quando modalidades fornecem informação contraditória ou ambígua, o sistema aplica regras de priorização hierárquicas:

\begin{enumerate}

    \item \textbf{Confiança (Confidence Score)}: Modalidade com maior confidence score prevalece \\
    Exemplo: Voz (0.92) vs Gesto (0.73) $\rightarrow$ Privilegia voz

    \item \textbf{Temporalidade}: Evento mais recente sobrepõe-se a antigos dentro da janela \\
    Exemplo: "zoom in" seguido de "zoom out" (< 2s) $\rightarrow$ Executa apenas "zoom out"

    \item \textbf{Tipo de Ação}: Certas ações privilegiam certas modalidades \\
    Exemplo: Pesquisa textual privilegia voz; Navegação espacial privilegia gestos

    \item \textbf{Contexto de Estado}: Estado atual do mapa influencia decisões \\
    Exemplo: Se já em navegação, comandos de navegação têm prioridade
\end{enumerate}

\section{Máquina de Estados SCXML}

O FusionEngine é implementado como uma máquina de estados finitos definida em SCXML, permitindo comportamento determinístico e extensível.

\subsection{Estados Principais}

\begin{description}
    \item[IDLE] Estado inicial de espera. O sistema aguarda eventos de modalidades. \\
    \textbf{Transições}: Receção de evento $\rightarrow$ COLLECTING

    \item[COLLECTING] Recolha de eventos dentro da janela temporal. Timer ativo de 2 segundos. \\
    \textbf{Transições}: Timeout ou evento de conclusão $\rightarrow$ FUSING

    \item[FUSING] Processamento e fusão dos eventos recolhidos. Aplicação de regras de fusão. \\
    \textbf{Transições}: Fusão bem-sucedida $\rightarrow$ VALIDATING; Conflito $\rightarrow$ CONFLICT\_RESOLUTION

    \item[CONFLICT\_RESOLUTION] Resolução de conflitos entre modalidades usando regras de priorização. \\
    \textbf{Transições}: Conflito resolvido $\rightarrow$ VALIDATING; Falha $\rightarrow$ ERROR

    \item[VALIDATING] Validação semântica e sintática do comando fundido. \\
    \textbf{Transições}: Válido $\rightarrow$ ROUTING; Inválido $\rightarrow$ ERROR

    \item[ROUTING] Encaminhamento do comando para o componente de execução apropriado. \\
    \textbf{Transições}: Enviado com sucesso $\rightarrow$ EXECUTING

    \item[EXECUTING] Aguarda confirmação de execução do Assistant. \\
    \textbf{Transições}: Execução concluída $\rightarrow$ FEEDBACK; Erro $\rightarrow$ ERROR

    \item[FEEDBACK] Geração e envio de feedback TTS ao utilizador. \\
    \textbf{Transições}: Feedback enviado $\rightarrow$ IDLE

    \item[ERROR] Tratamento de erros e geração de mensagem de erro. \\
    \textbf{Transições}: Sempre $\rightarrow$ IDLE (após notificar utilizador)
\end{description}

\subsection{Eventos e Transições}

\subsubsection{Eventos de Entrada (Input Events)}

\begin{itemize}
    \item \texttt{mmi.SPEECHIN}: Evento de reconhecimento de voz (origem: WebApp $\rightarrow$ RASA)
    \item \texttt{mmi.GESTUREIN}: Evento de reconhecimento de gesto (origem: Kinect $\rightarrow$ GenericGesturesModality)
    \item \texttt{mmi.STARTRESPONSE}: Confirmação de início de processamento
    \item \texttt{mmi.CANCELREQUEST}: Pedido de cancelamento de operação em curso
\end{itemize}

\subsubsection{Eventos Internos (Internal Events)}

\begin{itemize}
    \item \texttt{fusion.timeout}: Expiração da janela temporal de recolha
    \item \texttt{fusion.complete}: Fusão concluída com sucesso
    \item \texttt{fusion.conflict}: Conflito detetado entre modalidades
    \item \texttt{fusion.resolved}: Conflito resolvido com sucesso
    \item \texttt{validation.passed}: Comando validado
    \item \texttt{validation.failed}: Validação falhou
\end{itemize}

\subsubsection{Eventos de Saída (Output Events)}

\begin{itemize}
    \item \texttt{command.execute}: Comando pronto para execução (enviado ao Assistant via MMI Server)
    \item \texttt{tts.speak}: Mensagem de feedback TTS
    \item \texttt{error.notify}: Notificação de erro ao utilizador
\end{itemize}

\section{Processo de Fusão Detalhado}

\subsection{Fase 1: Receção e Normalização}

Quando um evento chega ao FusionEngine:

\begin{enumerate}
    \item \textbf{Parsing XML/JSON}: Evento MMI é parseado e validado
    \item \textbf{Extração de Dados}: Intent, entities, confidence e metadata são extraídos
    \item \textbf{Normalização}: Dados são normalizados para formato interno consistente
    \item \textbf{Timestamping}: Timestamp preciso é atribuído ao evento
    \item \textbf{Buffer Storage}: Evento é armazenado no buffer de fusão
\end{enumerate}

\subsection{Fase 2: Análise Semântica}

Para cada evento no buffer:

\begin{enumerate}
    \item \textbf{Classificação Semântica}: Evento é classificado como:
    \begin{itemize}
        \item Ação (e.g., zoom, pesquisar, navegar)
        \item Alvo (e.g., localização, estabelecimento)
        \item Modificador (e.g., modo de transporte, nível de zoom)
        \item Confirmação (afirmação/negação)
    \end{itemize}

    \item \textbf{Análise de Dependências}: Identifica relações entre eventos:
    \begin{itemize}
        \item Complementaridade (fornecem informação diferente)
        \item Redundância (fornecem informação igual)
        \item Conflito (fornecem informação contraditória)
    \end{itemize}

    \item \textbf{Construção de Frame}: Eventos são mapeados para slots de frame semântico
\end{enumerate}

\subsection{Fase 3: Fusão e Construção de Comando}

O algoritmo de fusão processa o frame semântico:

\begin{enumerate}
    \item \textbf{Merge de Slots}: Slots de diferentes eventos são combinados:
    \begin{itemize}
        \item Slots disjuntos são unidos diretamente
        \item Slots redundantes: seleciona maior confidence
        \item Slots conflituosos: aplica regras de priorização
    \end{itemize}

    \item \textbf{Inferência de Informação Implícita}:
    \begin{itemize}
        \item Uso de contexto do mapa (localização atual, zoom, resultados de pesquisa)
        \item Histórico de comandos recentes
        \item Preferências inferidas do utilizador
    \end{itemize}

    \item \textbf{Geração de Comando Unificado}: Criação de comando executável único
\end{enumerate}

\subsection{Fase 4: Validação}

O comando fundido é submetido a validação rigorosa:

\begin{enumerate}
    \item \textbf{Validação Sintática}: Verifica estrutura e tipos de dados
    \item \textbf{Validação Semântica}: Verifica coerência semântica (e.g., não é possível "navegar para zoom")
    \item \textbf{Validação de Precondições}: Verifica se estado atual permite execução (e.g., não é possível sair de Street View se não estiver em Street View)
    \item \textbf{Validação de Confidence}: Verifica se confidence agregado ultrapassa limiar mínimo
\end{enumerate}

\subsection{Fase 5: Routing e Execução}

Comando validado é enviado para execução:

\begin{enumerate}
    \item \textbf{Serialização}: Comando é serializado em formato MMI XML
    \item \textbf{Envio via MMI}: Publicação em MMI Server (porta 8005)
    \item \textbf{Receção pelo Assistant}: Assistant conectado via WebSocket recebe comando
    \item \textbf{Execução}: Assistant executa através de Selenium WebDriver
    \item \textbf{Confirmação}: Assistant envia feedback de sucesso/erro
\end{enumerate}

\section{Exemplos de Fusão}

\subsection{Exemplo 1: Fusão Voz + Gesto (Complementar)}

\textbf{Input}: Voz "procurar restaurantes" + Gesto swipe direita

\begin{enumerate}
    \item Eventos chegam ao FusionEngine:
    \begin{itemize}
        \item t=0ms: SPEECH intent=\texttt{search\_location}, entity=\texttt{restaurantes}
        \item t=450ms: GESTURE type=\texttt{swipe\_right}
    \end{itemize}

    \item FusionEngine identifica complementaridade:
    \begin{itemize}
        \item Voz fornece: tipo de POI (restaurantes)
        \item Gesto fornece: ação de navegação em lista
    \end{itemize}

    \item Decisão de fusão:
    \begin{itemize}
        \item Executar pesquisa de voz primeiro
        \item Após resultados aparecerem, executar navegação em lista
    \end{itemize}

    \item Comando fundido:
    \begin{verbatim}
{
  "sequence": [
    {"action": "search", "query": "restaurantes"},
    {"action": "navigate_results", "direction": "next"}
  ]
}
    \end{verbatim}
\end{enumerate}

\subsection{Exemplo 2: Fusão com Resolução de Conflito}

\textbf{Input}: Voz "afasta o zoom" + Gesto zoom in (mãos aproximam)

\begin{enumerate}
    \item FusionEngine deteta conflito:
    \begin{itemize}
        \item SPEECH: intent=\texttt{zoom\_out}, confidence=0.91
        \item GESTURE: type=\texttt{zoom\_in}, confidence=0.68
    \end{itemize}

    \item Aplicação de regras de priorização:
    \begin{itemize}
        \item Voz tem maior confidence (0.91 vs 0.68)
        \item Voz é mais explícita para ações de zoom
        \item \textbf{Decisão}: Privilegiar voz
    \end{itemize}

    \item Comando final: \texttt{zoom\_out} (voz prevalece)

    \item Logging do conflito para análise posterior
\end{enumerate}
\end{comment}


% ============================================================================
% Chapter 5: Modalidades de Interação
% ============================================================================

\chapter{Modalidades de Interação}

O sistema integra múltiplas modalidades de input e output que operam em conjunto através do motor de fusão. Este capítulo detalha a implementação técnica de cada modalidade e respetiva integração no sistema global.

\section{Modalidade de Voz com RASA NLU}

A modalidade de voz constitui a forma principal de interação, permitindo controlo através de linguagem natural em português europeu.

\begin{comment}
    \subsection{Arquitetura de Processamento de Voz}

O processamento de voz envolve três componentes principais:

\begin{enumerate}
    \item \textbf{Captura de Voz}: Web Speech API no browser captura áudio e realiza Speech-to-Text
    \item \textbf{Processamento NLU}: RASA processa texto e extrai intent + entities
    \item \textbf{Envio para Fusão}: WebApp envia evento MMI para FusionEngine
\end{enumerate}
\end{comment}

\begin{comment}
    \subsection{Pipeline RASA}

O sistema utiliza RASA 3.5 com o seguinte pipeline de processamento:

\begin{verbatim}
pipeline:
  - name: WhitespaceTokenizer
  - name: RegexFeaturizer
  - name: LexicalSyntacticFeaturizer
  - name: CountVectorsFeaturizer
  - name: CountVectorsFeaturizer
    analyzer: char_wb
    min_ngram: 1
    max_ngram: 4
  - name: DIETClassifier
    epochs: 100
    constrain_similarities: true
  - name: EntitySynonymMapper
  - name: ResponseSelector
    epochs: 100
\end{verbatim}

\subsubsection{Componentes do Pipeline}

\begin{description}
    \item[WhitespaceTokenizer] Divide texto em tokens usando espaços como delimitadores
    \item[RegexFeaturizer] Extrai features usando expressões regulares (e.g., padrões de números, datas)
    \item[LexicalSyntacticFeaturizer] Extrai features lexicais e sintáticas (POS tags)
    \item[CountVectorsFeaturizer] Gera vetores TF-IDF de palavras e caracteres (n-gramas)
    \item[DIETClassifier] Dual Intent Entity Transformer - classificador neural para intents e entities
    \item[EntitySynonymMapper] Mapeia sinónimos de entidades para valores canónicos
    \item[ResponseSelector] Seleciona resposta apropriada baseada no contexto
\end{description}
\end{comment}



\subsection{Intents Implementados}

O modelo RASA foi configurado para suportar um total de 26 \textit{intents} distintos. As tabelas que se seguem apresentam a especificação detalhada de cada \textbf{intent}, a sua função e ainda um exemplo representativo.


\begin{longtable}{p{0.30\textwidth} p{0.30\textwidth} p{0.33\textwidth}}
    \caption{Intent - Confirmação} 
    \label{tab:intents_confirmacao} \\
    
    % --- Cabeçalho da Primeira Página ---
    \toprule
    \textbf{Intent (ID)} & \textbf{Função} & \textbf{Exemplos de \textit{Utterances}} \\
    \midrule
    \endfirsthead
    \endlastfoot
    
    % 1. Categoria: Confirmação
    \texttt{affirm} & Afirmação positiva & "sim", "confirmo", "está bem" \\
    \texttt{deny} & Negação ou cancelamento & "não", "cancela", "abortar" \\
\end{longtable}

% 2. Categoria: Pesquisa e Navegação
\begin{longtable}{p{0.30\textwidth} p{0.30\textwidth} p{0.33\textwidth}}
    \caption{Intent - Pesquisa e Navegação} 
    \label{tab:intents_pesquisa_Navegacao} \\
    
    % --- Cabeçalho da Primeira Página ---
    \toprule
    \textbf{Intent (ID)} & \textbf{Função} & \textbf{Exemplos de \textit{Utterances}} \\
    \midrule
    \endfirsthead
    \endlastfoot
    
    \texttt{search\_location} & Pesquisar localização/POI & "procurar cafés", "ver Torre Eiffel" \\
    \texttt{get\_directions} & Obter rota entre pontos & "direções para Lisboa", "rota para Aveiro" \\
    \texttt{start\_navigation} & Iniciar navegação ativa & "iniciar navegação", "começar", "vai" \\
    \texttt{stop\_navigation} & Parar navegação ativa & "cancelar rota", "terminar" \\
\end{longtable}


% 3. Categoria: Controlo de Mapa
\begin{longtable}{p{0.30\textwidth} p{0.30\textwidth} p{0.33\textwidth}}
    \caption{Intent - Controlo de Mapa} 
    \label{tab:intents_controlo_mapa} \\
    
    % --- Cabeçalho da Primeira Página ---
    \toprule
    \textbf{Intent (ID)} & \textbf{Função} & \textbf{Exemplos} \\
    \midrule
    \endfirsthead
    \endlastfoot
    
    \texttt{zoom\_in} & Aumentar nível de zoom & "aproximar", "aumentar zoom" \\
    \texttt{zoom\_out} & Diminuir nível de zoom & "afastar", "diminuir zoom" \\
    \texttt{change\_map\_type} & Alterar camada visual & "mostrar satélite", "mudar para terreno" \\
    \texttt{recenter\_map} & Focar no utilizador & "centrar", "onde estou" \\
    \texttt{center\_location} & Focar local específico & -- \\
    \texttt{show\_traffic} / \texttt{hide\_traffic} & Alternar camada de tráfego & "mostrar trânsito", "esconder trânsito" \\
\end{longtable}


% 4. Categoria: Informação de Viagem
\begin{longtable}{p{0.30\textwidth} p{0.30\textwidth} p{0.33\textwidth}}
    \caption{Intent - Informação de Viagem} 
    \label{tab:intents_informacao_viagem} \\
    
    % --- Cabeçalho da Primeira Página ---
    \toprule
    \textbf{Intent (ID)} & \textbf{Função} & \textbf{Exemplos} \\
    \midrule
    \endfirsthead
    \endlastfoot
    
    \texttt{get\_trip\_duration} & Consultar tempo estimado & -- \\
    \texttt{get\_trip\_distance} & Consultar distância total & -- \\
    \texttt{change\_transport\_mode} & Mudar modo (carro/pé) & -- \\
    \texttt{swap\_route} & Escolher rota alternativa & -- \\
\end{longtable}


% 5. Categoria: Seleção
\begin{longtable}{p{0.30\textwidth} p{0.30\textwidth} p{0.33\textwidth}}
    \caption{Intent - Seleção} 
    \label{tab:intents_selecao} \\
    
    % --- Cabeçalho da Primeira Página ---
    \toprule
    \textbf{Intent (ID)} & \textbf{Função} & \textbf{Exemplos} \\
    \midrule
    \endfirsthead
    \endlastfoot

    \texttt{select\_place} & Selecionar item da lista & -- \\
    \texttt{select\_alternative\_route} & Confirmar rota opcional & -- \\
\end{longtable}


% 6. Categoria: Informação de Locais
\begin{longtable}{p{0.30\textwidth} p{0.30\textwidth} p{0.33\textwidth}}
    \caption{Intent - Informação de Locais} 
    \label{tab:intents_informacao_locais} \\
    
    % --- Cabeçalho da Primeira Página ---
    \toprule
    \textbf{Intent (ID)} & \textbf{Função} & \textbf{Exemplos} \\
    \midrule
    \endfirsthead
    \endlastfoot

    \texttt{show\_place\_details} & Expandir detalhes & -- \\
    \texttt{show\_reviews} & Ler avaliações & -- \\
    \texttt{show\_photos} & Visualizar galeria & -- \\
    \texttt{get\_opening\_hours} & Verificar horários & -- \\
    \texttt{get\_location\_info} & Obter informação do local & -- \\
    \texttt{whats\_here} & Explorar redondezas & -- \\
\end{longtable}


% 7. Categoria: Interação Geral
\begin{longtable}{p{0.30\textwidth} p{0.30\textwidth} p{0.33\textwidth}}
    \caption{Intent - Interação Geral} 
    \label{tab:intents_interacao_geral} \\
    
    % --- Cabeçalho da Primeira Página ---
    \toprule
    \textbf{Intent (ID)} & \textbf{Função} & \textbf{Exemplos} \\
    \midrule
    \endfirsthead
    \endlastfoot

    \texttt{greet} & Iniciar interação & "olá", "bom dia" \\
    \texttt{goodbye} & Encerrar interação & "adeus", "até logo" \\
    \texttt{thanks} & Agradecimento & "obrigado" \\
    \texttt{cancel} & Cancelar ação & "cancelar", "esquece" \\
    \texttt{help} & Solicitar assistência & "ajuda", "o que sabes fazer?" \\
\end{longtable}




\begin{comment}
    \subsubsection{Confirmação (2 intents)}
\begin{itemize}
    \item \texttt{affirm}: Afirmação (sim, confirmo, está bem)
    \item \texttt{deny}: Negação (não, cancela, abortar)
\end{itemize}

\subsubsection{Pesquisa e Navegação (4 intents)}
\begin{itemize}
    \item \texttt{search\_location}: Pesquisar localização ou POI \\
    Exemplos: "procurar cafés", "encontrar hospital mais próximo", "ver Torre Eiffel"

    \item \texttt{get\_directions}: Obter direções entre dois pontos \\
    Exemplos: "direções para Lisboa", "como chego ao Porto de carro", "rota para Aveiro"

    \item \texttt{start\_navigation}: Iniciar navegação turn-by-turn \\
    Exemplos: "iniciar navegação", "começar", "vai"

    \item \texttt{stop\_navigation}: Parar navegação \\
    Exemplos: "parar navegação", "cancelar rota", "terminar"
\end{itemize}

\subsubsection{Controlo de Mapa (6 intents)}
\begin{itemize}
    \item \texttt{zoom\_in}: Aumentar zoom \\
    Exemplos: "aproximar", "aumentar zoom", "zoom in"

    \item \texttt{zoom\_out}: Diminuir zoom \\
    Exemplos: "afastar", "diminuir zoom", "zoom out"

    \item \texttt{change\_map\_type}: Alterar tipo de vista \\
    Exemplos: "mostrar satélite", "mudar para terreno", "vista normal"

    \item \texttt{recenter\_map}: Recentrar na localização atual \\
    Exemplos: "centrar", "onde estou", "mostrar a minha localização"

    \item \texttt{center\_location}: Centrar numa localização específica

    \item \texttt{show\_traffic / hide\_traffic}: Mostrar/ocultar tráfego
\end{itemize}

\subsubsection{Informação de Viagem (4 intents)}
\begin{itemize}
    \item \texttt{get\_trip\_duration}: Consultar duração estimada
    \item \texttt{get\_trip\_distance}: Consultar distância total
    \item \texttt{change\_transport\_mode}: Alterar modo de transporte
    \item \texttt{swap\_route}: Trocar para rota alternativa
\end{itemize}

\subsubsection{Seleção (2 intents)}
\begin{itemize}
    \item \texttt{select\_place}: Selecionar lugar da lista
    \item \texttt{select\_alternative\_route}: Selecionar rota alternativa
\end{itemize}

\subsubsection{Informação de Locais (5 intents)}
\begin{itemize}
    \item \texttt{show\_place\_details}: Mostrar detalhes de lugar
    \item \texttt{show\_reviews}: Mostrar avaliações
    \item \texttt{show\_photos}: Mostrar fotografias
    \item \texttt{get\_opening\_hours}: Consultar horários
    \item \texttt{whats\_here}: Descobrir o que há numa localização
\end{itemize}

\subsubsection{Interação Geral (4 intents)}
\begin{itemize}
    \item \texttt{greet}: Saudação
    \item \texttt{goodbye}: Despedida
    \item \texttt{thanks}: Agradecimento
    \item \texttt{help}: Pedido de ajuda
\end{itemize}
\end{comment}

\subsection{Entidades Extraídas}

Complementarmente à classificação de intenções, o sistema reconhece 6 tipos de entidades principais. A Tabela \ref{tab:entidades_rasa} detalha as entidades modeladas, os seus domínios e exemplos.

\begin{longtable}{p{0.18\textwidth} p{0.37\textwidth} p{0.38\textwidth}}
\caption{Lista de entidades}
\label{tab:entidades_rasa} \\

% --- Cabeçalho ---
\toprule
\textbf{Entidade} & \textbf{Descrição e Domínio} & \textbf{Exemplo de Extração} \\
\midrule
\endfirsthead

% --- Cabeçalho (Continuação) ---
\multicolumn{3}{c}%
{{\bfseries \tablename\ \thetable{} -- Continuação da página anterior}} \\
\toprule
\textbf{Entidade} & \textbf{Descrição e Domínio} & \textbf{Exemplo de Extração} \\
\midrule
\endhead

% --- Rodapé ---
\bottomrule
\endlastfoot

% 1. Location
\texttt{location} & Identifica topónimos, moradas genéricas ou Pontos de Interesse (POIs). & "procurar \textbf{[cafés em Lisboa]}(location)" \\
\addlinespace

% 2. Destination
\texttt{destination} & O ponto de chegada num pedido de navegação ou cálculo de rota. & "direções para \textbf{[Porto]}(destination)" \\
\addlinespace

% 3. Origin
\texttt{origin} & O ponto de partida (quando diferente da localização atual do utilizador). & "de \textbf{[Aveiro]}(origin) para..." \\
\addlinespace

% 4. Transport Mode
\texttt{transport\_mode} & Modalidade de deslocação pretendida. \newline \textit{\small Valores: carro, transportes públicos, a pé, bicicleta} & "rota de \textbf{[carro]}(transport\_mode)" \\
\addlinespace

% 5. Map Type
\texttt{map\_type} & Estilo visual da camada do mapa. \newline \textit{\small Valores: satélite, terreno, mapa, padrão} & "vista de \textbf{[satélite]}(map\_type)" \\
\addlinespace

% 6. Zoom Level
\texttt{zoom\_level} & Qualificador de intensidade para operações de zoom. \newline \textit{\small Valores: muito, pouco} & "aproximar \textbf{[muito]}(zoom\_level)" \\

\end{longtable}

\begin{comment}
    \subsection{Treino do Modelo}

O modelo foi treinado com corpus extensivo de exemplos em português europeu:

\begin{itemize}
    \item \textbf{Exemplos de treino}: Mais de 300 enunciados anotados
    \item \textbf{Variações linguísticas}: Múltiplas formas de expressar cada intent
    \item \textbf{Entidades anotadas}: Markup manual de todas as entidades
    \item \textbf{Sinónimos}: Mapeamento de variações para valores canónicos
    \item \textbf{Epochs}: 100 épocas de treino para DIETClassifier
    \item \textbf{Validação}: Split treino/teste 80/20 com avaliação de precisão
\end{itemize}
\end{comment}


\subsection{Integração com FusionEngine}

O fluxo de processamento de voz é:

\begin{enumerate}
    \item Utilizador fala para o microfone
    \item Web Speech API captura áudio e converte em texto (browser-side)
    \item WebApp envia texto via HTTP POST para RASA Server (porta 5005):
    \begin{comment}
        \begin{verbatim}
POST http://localhost:5005/model/parse
{"text": "procurar restaurantes em Aveiro"}
    \end{verbatim}
    \end{comment}

    \item RASA processa e retorna JSON com intent + entities:
    \begin{comment}
        \begin{verbatim}
{
  "intent": {"name": "search_location", "confidence": 0.94},
  "entities": [{"entity": "location", "value": "restaurantes em Aveiro"}],
  "text": "procurar restaurantes em Aveiro"
}
    \end{verbatim}
    \end{comment}

    \item WebApp constrói evento MMI e envia para FusionEngine (porta 8001):
    \begin{comment}
        \begin{verbatim}
POST https://127.0.0.1:8001/IM/USER1/SPEECHIN
<mmi:mmi>
  <mmi:startRequest>
    <mmi:data>
      <command>
        {
          "recognized": ["SPEECH"],
          "nlu": {...}
        }
      </command>
    </mmi:data>
  </mmi:startRequest>
</mmi:mmi>
    \end{verbatim}
    \end{comment}

    \item FusionEngine processa evento conforme descrito no Capítulo 4
\end{enumerate}

\section{Modalidade de Gestos com Kinect v2}

A modalidade de gestos permite interação espacial intuitiva através do sensor Kinect v2.

\begin{comment}
    \subsection{Arquitetura de Reconhecimento de Gestos}

\subsubsection{Hardware: Sensor Kinect v2}

O Kinect v2 fornece:
\begin{itemize}
    \item \textbf{Câmara RGB}: 1920x1080 @ 30fps
    \item \textbf{Sensor de Profundidade}: 512x424 time-of-flight
    \item \textbf{Tracking de Skeleton}: 25 pontos articulares por pessoa, até 6 pessoas
    \item \textbf{Hand State Detection}: Deteção de estado da mão (aberta, fechada, apontando)
\end{itemize}

\subsubsection{Software: GenericGesturesModality}

Componente Java que:
\begin{enumerate}
    \item Recebe dados skeleton do Kinect SDK
    \item Processa coordenadas 3D das articulações
    \item Aplica algoritmos de deteção de gestos
    \item Gera eventos MMI para o FusionEngine
\end{enumerate}

\newpage
\end{comment}

\subsection{Gestos Reconhecidos}

O sistema implementa num total 15 gestos distintos organizados por funcionalidade referentes à modalidade de gestos

\begin{longtable}{p{0.25\textwidth} p{0.30\textwidth} p{0.35\textwidth}}
    \caption{Lista de gestos - Navegação Espacial} 
    \label{tab:navegacao-espacial} \\

    % Cabeçalho da Tabela
    \toprule
    \textbf{Gesto} & \textbf{Ação / Significado} & \textbf{Detalhes do Algoritmo} \\
    \midrule
    \endfirsthead
    \endlastfoot

    % --- Navegação Espacial ---
    \texttt{SwipeUp} & Move mapa para norte & Vel. vertical > limiar + deslocamento > 20cm \\
    \texttt{SwipeDown} & Move mapa para sul & Vel. vertical negativa > limiar \\
    \texttt{SwipeLeft\_Left} & Move mapa para oeste & Desloc. horizontal (mão esq.) < -20cm em < 1s \\
    \texttt{SwipeRight\_Right} & Move mapa para este & Desloc. horizontal (mão dir.) > 20cm em < 1s \\
\end{longtable}

\begin{longtable}{p{0.25\textwidth} p{0.30\textwidth} p{0.35\textwidth}}
    \caption{Lista de gestos - Gestos de Zoom}
    \label{tab:gestos-zoom} \\
    
    % Cabeçalho da Tabela
    \toprule
    \textbf{Gesto} & \textbf{Ação / Significado} & \textbf{Detalhes do Algoritmo} \\
    \midrule
    \endfirsthead
    \endlastfoot

    % --- Zoom ---
    \texttt{ZoomIn} & Aumentar zoom & Distância entre mãos diminui > 15cm \\
    \texttt{ZoomOut} & Diminuir zoom & Distância entre mãos aumenta > 15cm \\
\end{longtable}

\begin{longtable}{p{0.25\textwidth} p{0.30\textwidth} p{0.35\textwidth}}
    \caption{Lista de gestos - Gestos de Confirmação} 
    \label{tab:gestos-confirmacao} \\
    
    % Cabeçalho da Tabela
    \toprule
    \textbf{Gesto} & \textbf{Ação / Significado} & \textbf{Detalhes do Algoritmo} \\
    \midrule
    \endfirsthead
    \endlastfoot
    
    % --- Confirmação ---
    \texttt{ThumbsUp} (Yes/OK) & Afirmação & Mão fechada + polegar cima + acima do ombro \\
\end{longtable}



\begin{longtable}{p{0.25\textwidth} p{0.30\textwidth} p{0.35\textwidth}}
    \caption{Lista de gestos - Funcionalidades} 
    \label{tab:gestos_funcionalidades} \\
    
    % Cabeçalho da Tabela
    \toprule
    \textbf{Gesto} & \textbf{Ação / Significado} & \textbf{Detalhes do Algoritmo} \\
    \midrule
    \endfirsthead
    \endlastfoot
    
    % --- Funcionalidades ---
    \texttt{Hotels} & Filtro "Hotéis" & Ativa filtro no Google Maps \\
    \texttt{Restaurants} & Filtro "Restaurantes" & Ativa filtro de restaurantes \\
    \texttt{Transports} & Filtro "Transportes" & Ativa filtro de transportes públicos \\
    \texttt{Gas Stations} & Filtro "Bombas de Gasolina" & Ativa filtro de bombas de gasolina \\
    \texttt{Camera} & Coisas a fazer & Abre painel de pontos turísticos \\
\end{longtable}


\begin{longtable}{p{0.25\textwidth} p{0.30\textwidth} p{0.35\textwidth}}
    \caption{Lista de gestos - Street View} 
    \label{tab:gestos-street-view} \\
    
    % Cabeçalho da Tabela
    \toprule
    \textbf{Gesto} & \textbf{Ação / Significado} & \textbf{Detalhes do Algoritmo} \\
    \midrule
    \endfirsthead
    \endlastfoot
    
    % --- Street View ---
    \texttt{EnterStreet} & Entrar no modo Street View & Gesto de "push" (empurrar) com ambas as mãos \\
    \texttt{ExitStreet} & Sair do modo Street View & Gesto de "pull" (puxar) com ambas as mãos \\
    \texttt{Forward} & Avançar no Street View & Gesto de caminhar/avançar \\
\end{longtable}

\begin{longtable}{p{0.25\textwidth} p{0.30\textwidth} p{0.35\textwidth}}
    \caption{Lista de gestos - Seleção e Navegação} 
    \label{tab:gestos_selecao_navegacao} \\
    
    % Cabeçalho da Tabela
    \toprule
    \textbf{Gesto} & \textbf{Ação / Significado} & \textbf{Detalhes do Algoritmo} \\
    \midrule
    \endfirsthead
    \endlastfoot
    
    % --- Seleção e Navegação ---
    \texttt{Select} & Selecionar elemento & Mão em estado "apontar" estável por > 1s \\
    \texttt{UpOption\_Right} & Subir na lista & Swipe up (mão direita) em contexto de lista \\
    \texttt{DownOption\_Left} & Descer na lista & Swipe down (mão esquerda) em contexto de lista \\
\end{longtable}


\begin{comment}
    

O sistema implementa 15 gestos distintos organizados por funcionalidade:

\subsubsection{Gestos de Navegação Espacial}

\begin{description}
    \item[SwipeUp] Deslizar mão para cima - Move mapa para norte \\
    \textbf{Algoritmo}: Velocidade vertical da mão > limiar + deslocamento > 20cm

    \item[SwipeDown] Deslizar mão para baixo - Move mapa para sul \\
    \textbf{Algoritmo}: Velocidade vertical negativa da mão > limiar

    \item[SwipeLeft\_Left] Deslizar mão esquerda para esquerda - Move mapa para oeste \\
    \textbf{Algoritmo}: Deslocamento horizontal mão esquerda < -20cm em < 1s

    \item[SwipeRight\_Right] Deslizar mão direita para direita - Move mapa para este \\
    \textbf{Algoritmo}: Deslocamento horizontal mão direita > 20cm em < 1s
\end{description}

\subsubsection{Gestos de Zoom}

\begin{description}
    \item[ZoomIn] Aproximar mãos - Aumentar zoom \\
    \textbf{Algoritmo}: Distância entre mãos diminui > 15cm em janela temporal

    \item[ZoomOut] Afastar mãos - Diminuir zoom \\
    \textbf{Algoritmo}: Distância entre mãos aumenta > 15cm em janela temporal
\end{description}

\subsubsection{Gestos de Confirmação}

\begin{description}
    \item[ThumbsUp / ThumbUp / Yes / OK] Polegar para cima - Afirmação \\
    \textbf{Algoritmo}: Mão fechada + polegar estendido vertical + posição acima do ombro

    \item[ThumbsDown / ThumbDown / No / Cancel] Polegar para baixo - Negação \\
    \textbf{Algoritmo}: Mão fechada + polegar estendido vertical + posição abaixo do ombro
\end{description}

\subsubsection{Gestos de Funcionalidades}

\begin{description}
    \item[Hotels] Gesto específico para filtro de hotéis \\
    \textbf{Mapeamento}: Ativa filtro "Hotéis" no Google Maps

    \item[Restaurants] Gesto para filtro de restaurantes \\
    \textbf{Mapeamento}: Ativa filtro "Restaurantes"

    \item[Transports] Gesto para filtro de transportes \\
    \textbf{Mapeamento}: Ativa filtro "Transportes públicos"

    \item[Camera] Gesto para "Coisas a fazer" \\
    \textbf{Mapeamento}: Abre painel de pontos turísticos
\end{description}

\subsubsection{Gestos de Street View}

\begin{description}
    \item[EnterStreet] Entrar em modo Street View \\
    \textbf{Algoritmo}: Gesto de "push" (empurrar) com ambas as mãos

    \item[ExitStreet] Sair de modo Street View \\
    \textbf{Algoritmo}: Gesto de "pull" (puxar) com ambas as mãos
\end{description}

\subsubsection{Gestos de Seleção e Navegação}

\begin{description}
    \item[Select] Selecionar elemento \\
    \textbf{Algoritmo}: Mão em estado "apontar" estável por > 1s

    \item[UpOption\_Right] Subir em lista de opções \\
    \textbf{Algoritmo}: Swipe up com mão direita em contexto de lista

    \item[DownOption\_Left] Descer em lista de opções \\
    \textbf{Algoritmo}: Swipe down com mão esquerda em contexto de lista
\end{description}
\end{comment}


\subsection{Integração com FusionEngine}

Quando gesto é reconhecido:

\begin{enumerate}
    \item GenericGesturesModality gera evento MMI:
    \begin{comment}
        \begin{verbatim}
POST https://127.0.0.1:8001/IM/USER1/GESTUREIN
<mmi:mmi>
  <mmi:startRequest>
    <mmi:data>
      <command>
        {
          "recognized": ["GESTURES", "ZoomI"],
          "confidence": 0.84
        }
      </command>
    </mmi:data>
  </mmi:startRequest>
</mmi:mmi>
    \end{verbatim}
    \end{comment}

    \item FusionEngine recebe evento e extrai informação

    \item No Assistant, código semântico é mapeado para intent:
\begin{comment}
        \begin{verbatim}
gesture_to_intent = {
    "zoomi": "gesture_zoom_in",
    "zoomo": "gesture_zoom_out",
    "swipeu": "gesture_swipe_up",
    ...
}
    \end{verbatim}
\end{comment}

    \item Intent é processado pelo handler correspondente no Assistant
\end{enumerate}

\section{Modalidade de Saída: Text-to-Speech}

O sistema fornece feedback áudio através de síntese de voz em português. A tabela \ref{tab:feedback_tts} mostra diferentes tipos de feedback TTS.

\begin{comment}
    \subsection{Tecnologia: Web Speech API}

A implementação utiliza Web Speech API (\texttt{speechSynthesis}) disponível em browsers modernos:

\begin{verbatim}
const utterance = new SpeechSynthesisUtterance(text);
utterance.lang = 'pt-PT';
utterance.rate = 1.0;  // Velocidade normal
utterance.pitch = 1.0; // Tom neutro
utterance.volume = 1.0; // Volume máximo
speechSynthesis.speak(utterance);
\end{verbatim}

\subsubsection{Voz Utilizada}

\begin{itemize}
    \item \textbf{Nome}: Microsoft Maria Desktop (Windows) ou equivalente do sistema
    \item \textbf{Língua}: Português de Portugal (pt-PT)
    \item \textbf{Qualidade}: Voz neural de alta qualidade
    \item \textbf{Expressividade}: Entoação natural para português
\end{itemize}

\subsection{Arquitetura de Distribuição TTS}

O sistema implementa arquitetura publish-subscribe para TTS:

\begin{enumerate}
    \item \textbf{Assistant} (Publisher):
    \begin{itemize}
        \item Gera mensagem de feedback após executar comando
        \item Conecta via WebSocket ao TTS Server (porta 8083)
        \item Envia mensagem em formato texto JSON
    \end{itemize}

    \item \textbf{TTS WebSocket Server} (Broker):
    \begin{verbatim}
async def tts_handler(websocket):
    connected_clients.add(websocket)
    async for message in websocket:
        # Broadcast to all clients except sender
        for client in connected_clients:
            if client != websocket:
                await client.send(message)
    \end{verbatim}

    \item \textbf{WebApp} (Subscriber):
    \begin{itemize}
        \item Mantém WebSocket conectado ao TTS Server
        \item Recebe mensagens de feedback
        \item Executa síntese de voz via Web Speech API
    \end{itemize}
\end{enumerate}
\end{comment}

\begin{longtable}{p{0.25\textwidth} p{0.25\textwidth} p{0.40\textwidth}}
    \caption{Tipos de Feedback TTS e Exemplos de Mensagens} 
    \label{tab:feedback_tts} \\
    
    % Cabeçalho da Tabela
    \toprule
    \textbf{Tipo de Feedback} & \textbf{Descrição} & \textbf{Exemplos} \\
    \midrule
    \endfirsthead
    \endlastfoot
    
    % --- Conteúdo ---
    
    \textbf{Confirmação de Ações} & 
    Confirmação de execução bem-sucedida & 
    "Pesquisei por restaurantes em Aveiro"; 
    "Aumentei o zoom no mapa"; 
    "A obter direções para Lisboa"; 
    "Centrei o mapa na tua localização" \\
    \midrule
    
    \textbf{Reportagem de Resultados} & 
    Mensagens informativas sobre resultados obtidos & 
    "Encontrei 15 restaurantes na área"; 
    "A duração estimada é de 2h e 30m"; 
    "A distância total é de 180 km";
    "Existem 3 rotas alternativas" \\
    \midrule
    
    \textbf{Solicitação de Confirmação} & 
    Pedido de confirmação antes de ações críticas & 
    "Queres iniciar a navegação para Porto?"; 
    "Confirmas que queres mudar para modo de bicicleta?" \\
    \midrule
    
    \textbf{Notificação de Erros} & 
    Mensagens de erro claras e orientativas & 
    "Desculpa, não percebi. Podes repetir?"; 
    "Não encontrei resultados para essa pesquisa"; 
    "Ocorreu um erro ao carregar o mapa"; 
    "Precisas de especificar um destino" \\
    \midrule
    
    \textbf{Saudações e Interação} & 
    Mensagens de interação social & 
    "Boas! Eu sou a Assistente. Como te posso ajudar?"; 
    "De nada! Estou aqui se precisares"; 
    "Adeus! Até à próxima!" \\
    
    \bottomrule
\end{longtable}

\begin{comment}

\section{Módulo de Execução: Selenium WebDriver}

O Assistant é responsável por executar comandos no Google Maps através de automação web.
    
\subsection{Arquitetura Selenium}

\subsubsection{Configuração do WebDriver}

\begin{verbatim}
from selenium import webdriver
from selenium.webdriver.chrome.service import Service
from selenium.webdriver.chrome.options import Options

options = Options()
options.add_argument('--start-maximized')
options.add_argument('--disable-blink-features=AutomationControlled')
# Opcional: headless mode
# options.add_argument('--headless')

driver = webdriver.Chrome(options=options)
driver.get('https://www.google.com/maps')
\end{verbatim}

\subsubsection{Gestão de Sessão Persistente}

O driver mantém sessão persistente durante toda a execução:
\begin{itemize}
    \item Browser permanece aberto entre comandos
    \item Estado do mapa (zoom, posição, pesquisas) é preservado
    \item Cookies e autenticação mantidos
    \item Reduz latência (não reinicia browser a cada comando)
\end{itemize}

\subsection{Operações Implementadas}

\subsubsection{Pesquisa (search\_location)}

\begin{verbatim}
def search_location(driver, query):
    # Localizar caixa de pesquisa
    search_box = WebDriverWait(driver, 10).until(
        EC.presence_of_element_located((By.ID, "searchboxinput"))
    )

    # Limpar e preencher
    search_box.clear()
    search_box.send_keys(query)
    search_box.send_keys(Keys.RETURN)

    # Aguardar resultados carregarem
    WebDriverWait(driver, 10).until(
        EC.presence_of_element_located((By.CSS_SELECTOR, ".section-result"))
    )
\end{verbatim}

\subsubsection{Zoom (zoom\_in / zoom\_out)}

Implementação via JavaScript para precisão:
\begin{verbatim}
def zoom_in(driver, amount=1):
    for _ in range(amount):
        driver.execute_script("""
            var map = document.querySelector('[role="region"]');
            var zoomIn = document.querySelector('[aria-label*="Zoom in"]');
            if (zoomIn) zoomIn.click();
        """)
        time.sleep(0.1)  # Delay para animação
\end{verbatim}

Alternativamente, via API do Google Maps (se acessível):
\begin{verbatim}
driver.execute_script("""
    var currentZoom = map.getZoom();
    map.setZoom(currentZoom + 1);
""")
\end{verbatim}

\subsubsection{Obter Direções (get\_directions)}

\begin{verbatim}
def get_directions(driver, destination, origin=None, mode='driving'):
    # Construir URL com parâmetros
    base_url = "https://www.google.com/maps/dir/"

    if origin:
        url = f"{base_url}{origin}/{destination}/@?travelmode={mode}"
    else:
        url = f"{base_url}/{destination}/@?travelmode={mode}"

    driver.get(url)

    # Aguardar painel de direções
    WebDriverWait(driver, 15).until(
        EC.presence_of_element_located((By.ID, "section-directions-trip-0"))
    )
\end{verbatim}

Modos de transporte suportados:
\begin{itemize}
    \item \texttt{driving}: Carro
    \item \texttt{transit}: Transportes públicos
    \item \texttt{walking}: A pé
    \item \texttt{bicycling}: Bicicleta
\end{itemize}

\subsubsection{Alterar Tipo de Mapa (change\_map\_type)}

\begin{verbatim}
def change_map_type(driver, map_type):
    # Abrir menu de camadas
    layers_btn = driver.find_element(By.CSS_SELECTOR,
                                      '[aria-label*="Layers"]')
    layers_btn.click()

    # Selecionar tipo
    type_mapping = {
        "satellite": "Satellite",
        "terrain": "Terrain",
        "default": "Default map"
    }

    type_btn = driver.find_element(By.XPATH,
                                    f"//button[contains(., '{type_mapping[map_type]}')]")
    type_btn.click()
\end{verbatim}

\subsection{Tratamento de Erros e Robustez}

\subsubsection{Sistema Centralizado de Error Handling}

O sistema implementa um módulo centralizado de tratamento de erros (\texttt{error\_handler.py}) que converte exceções técnicas do Selenium em mensagens compreensíveis em português para o TTS.

\textbf{Arquitetura do ErrorHandler:}

\begin{enumerate}
    \item \textbf{Categorização Automática}: Exceções são automaticamente categorizadas em 7 tipos (rede, browser, elemento não encontrado, timeout, interação, validação, desconhecido)
    
    \item \textbf{Mensagens User-Friendly}: Cada categoria tem mensagens em português otimizadas para TTS:
    \begin{itemize}
        \item NoSuchElementException → "Não encontrei o elemento no mapa"
        \item TimeoutException → "A operação demorou muito tempo. Tenta outra vez"
        \item WebDriverException (network) → "Verifica a tua ligação à Internet"
    \end{itemize}
    
    \item \textbf{Contexto Operacional}: Funções auxiliares específicas para cada tipo de operação:
\end{enumerate}

\begin{verbatim}
# Exemplo de utilização nos intent handlers
try:
    result = maps_page.search_location(query)
except Exception as e:
    user_message = handle_search_error(e, query)
    # TTS diz: "Não consegui procurar por [query]"
    return IntentResponse(success=False, message=user_message)
\end{verbatim}

\textbf{Aplicação no Sistema:}
\begin{itemize}
    \item \textbf{35 exception handlers} atualizados em 7 ficheiros
    \item \textbf{Logging técnico}: Exceções completas registadas para debugging
    \item \textbf{Output TTS}: Apenas mensagens compreensíveis para o utilizador
    \item \textbf{Fallback messages}: Mensagens específicas por contexto (pesquisa, direções, navegação, controlo do mapa)
\end{itemize}

\textbf{Benefícios:}
\begin{enumerate}
    \item Melhor UX: Utilizador compreende o que correu mal sem jargão técnico
    \item Debugging facilitado: Logs técnicos completos mantidos para developers
    \item Manutenção centralizada: Mensagens geridas num único ponto
    \item Consistência: Todas as operações usam o mesmo formato de erro
\end{enumerate}

\subsubsection{Custom Exceptions}

Duas exceções personalizadas para casos específicos:

\begin{verbatim}
class ValidationError(Exception):
    """Erros de validação com mensagem user-friendly"""
    def __init__(self, user_message: str, technical_details: str):
        self.user_message = user_message  # Para TTS
        self.technical_details = technical_details  # Para logs

class OperationFailedError(Exception):
    """Operações falhadas com contexto"""
    def __init__(self, operation: str, reason: str):
        self.operation = operation
        self.reason = reason
\end{verbatim}

Estas exceções permitem separar claramente a mensagem para o utilizador (TTS) dos detalhes técnicos (logs), garantindo uma experiência robusta e profissional.


\subsubsection{Gestão de Timeouts}

\begin{itemize}
    \item \textbf{Implicit Wait}: 10 segundos (procura de elementos)
    \item \textbf{Page Load Timeout}: 30 segundos (carregamento de página)
    \item \textbf{Script Timeout}: 15 segundos (execução de JavaScript)
\end{itemize}

\begin{verbatim}
driver.implicitly_wait(10)
driver.set_page_load_timeout(30)
driver.set_script_timeout(15)
\end{verbatim}

\subsubsection{Recuperação de Crashes}

\begin{verbatim}
def ensure_browser_alive(driver):
    try:
        driver.current_url  # Test if browser is responsive
    except:
        logger.error("Browser crashed, restarting...")
        driver_manager.stop()
        driver = driver_manager.start()
        driver.get('https://www.google.com/maps')
    return driver
\end{verbatim}

\subsection{Extração de Informação}

\subsubsection{Obter Duração de Viagem}

\begin{verbatim}
def get_trip_duration(driver):
    duration_elem = driver.find_element(By.CSS_SELECTOR,
                                         '.section-directions-trip-duration')
    return duration_elem.text  # e.g., "2 h 30 min"
\end{verbatim}

\subsubsection{Obter Distância}

\begin{verbatim}
def get_trip_distance(driver):
    distance_elem = driver.find_element(By.CSS_SELECTOR,
                                         '.section-directions-trip-distance')
    return distance_elem.text  # e.g., "180 km"
\end{verbatim}

\subsubsection{Extrair Resultados de Pesquisa}

\begin{verbatim}
def get_search_results(driver):
    results = []
    result_elements = driver.find_elements(By.CSS_SELECTOR,
                                            '.section-result')

    for elem in result_elements[:5]:  # Top 5
        name = elem.find_element(By.CSS_SELECTOR, '.section-result-title').text
        rating = elem.find_element(By.CSS_SELECTOR, '.section-result-rating').text
        results.append({"name": name, "rating": rating})

    return results
\end{verbatim}
\end{comment}


% ============================================================================
% Chapter 6: Conclusões
% ============================================================================

\chapter{Conclusões}

Este capítulo apresenta uma síntese do trabalho realizado, avalia o cumprimento dos objetivos estabelecidos, discute as dificuldades encontradas e aponta direções para desenvolvimento futuro.

\section{Síntese do Trabalho Realizado}

O presente trabalho consistiu no desenvolvimento de um sistema multimodal para controlo de mapas digitais que integra reconhecimento de voz em português, reconhecimento de gestos através de Kinect v2, e um motor de fusão baseado em SCXML. O sistema implementado compreende cinco componentes principais operando de forma distribuída:

\begin{description}
    \item[FusionEngine] Motor de fusão multimodal baseado em máquina de estados SCXML, responsável por combinar informação de diferentes modalidades e gerar comandos unificados

    \item[MMI Framework] Implementação do protocolo W3C MMI que facilita comunicação padronizada entre modalidades e motor de fusão

    \item[RASA NLU Server] Sistema de processamento de linguagem natural treinado para reconhecer 26 intents em português europeu, com suporte para extração de 6 tipos de entidades

    \item[GenericGesturesModality] Módulo de reconhecimento de gestos que processa dados skeleton do Kinect v2 e identifica 15 gestos distintos

    \item[Google Maps Assistant] Módulo de execução que automatiza o Google Maps através de Selenium WebDriver e fornece feedback TTS

    \item[WebApp Server] Interface web que captura voz via Web Speech API e distribui mensagens TTS aos clientes
\end{description}

A arquitetura implementada segue princípios de modularidade e separação de responsabilidades, permitindo que cada componente opere de forma independente enquanto comunica através de protocolos bem definidos (MMI, WebSocket, HTTPS).

\section{Cumprimento de Objetivos}

O sistema desenvolvido cumpre os requisitos funcionais e não-funcionais estabelecidos no Capítulo 3:

\subsection{Objetivos Funcionais Alcançados}

\begin{itemize}
    \item \textbf{Reconhecimento de Voz}: Sistema reconhece comandos de voz em português europeu com taxa de precisão superior a 85\% para comandos bem formulados. O modelo RASA treinado suporta 26 intents distintos cobrindo pesquisa, navegação, controlo de mapa, e obtenção de informação.

    \item \textbf{Reconhecimento de Gestos}: Implementação de 15 gestos reconhecíveis pelo Kinect v2, incluindo navegação espacial (swipes), zoom, confirmação (thumbs up/down), e ativação de filtros. Taxa de reconhecimento superior a 80\% em condições adequadas de iluminação.

    \item \textbf{Fusão Multimodal}: Motor de fusão baseado em SCXML combina eventos de voz e gestos através de janela temporal de 2 segundos, com resolução de conflitos baseada em regras de priorização (confidence, especificidade, temporalidade).

    \item \textbf{Controlo Completo de Google Maps}: Automação via Selenium WebDriver permite execução de todas as operações principais: pesquisa, direções, navegação, zoom, alteração de vista, consulta de informação de locais.

    \item \textbf{Feedback TTS}: Sistema de Text-to-Speech em português fornece confirmação de ações, reportagem de resultados, notificação de erros, e orientação ao utilizador.

    \item \textbf{Performance}: Latências medidas dentro dos limites estabelecidos:
    \begin{itemize}
        \item Reconhecimento de voz: 1.2--1.8s (média 1.5s)
        \item Reconhecimento de gestos: 200--400ms (média 300ms)
        \item Fusão: 50--150ms (média 100ms)
        \item Execução de comandos: 0.5--2s dependendo da operação
    \end{itemize}
\end{itemize}

\subsection{Contribuições Técnicas}

O trabalho apresenta as seguintes contribuições:

\begin{enumerate}
    \item \textbf{Fusão Híbrida}: Implementação de fusão que combina abordagens baseadas em regras e fusão temporal, adaptada ao contexto de controlo de mapas

    \item \textbf{Corpus Português}: Criação de corpus de treino para RASA com mais de 300 enunciados em português europeu, cobrindo domínio de mapas e navegação

    \item \textbf{Mapeamento Semântico}: Sistema de mapeamento entre códigos semânticos de gestos Kinect e intents do sistema, permitindo fusão consistente

    \item \textbf{Arquitetura Distribuída}: Design de sistema multimodal distribuído usando WebSockets e protocolo MMI, facilmente extensível para novas modalidades

    \item \textbf{Robustez}: Implementação de mecanismos de tolerância a falhas (retry logic, recovery, timeouts) que garantem operação estável
\end{enumerate}

\section{Dificuldades Encontradas e Soluções}

Durante o desenvolvimento, foram enfrentados diversos desafios técnicos que exigiram investigação e adaptação:

\subsection{Compatibilidade de Versões Python}

\textbf{Problema}: RASA 3.6+ requer Python 3.10 ou 3.11, incompatível com Python 3.12 devido a dependências TensorFlow desatualizadas.

\textbf{Solução}: Criação de ambiente virtual com Python 3.10 especificamente para o projeto. Atualização de scripts de inicialização para ativar ambiente correto antes de lançar componentes Python.

\textbf{Lição}: Verificar matriz de compatibilidade de frameworks de machine learning antes de escolher versão de Python.

\subsection{Mudanças na API WebSockets}

\textbf{Problema}: Biblioteca \texttt{websockets} versão 13.x removeu parâmetro \texttt{path} do handler, causando \texttt{TypeError} em código legacy.

\textbf{Solução}: Refatoração de handlers WebSocket

\textbf{Lição}: Ler changelogs de bibliotecas ao atualizar versões, especialmente para mudanças breaking.

\subsection{Gestão de Certificados SSL}

\textbf{Problema}: Sistema requer certificados SSL em múltiplos componentes (FusionEngine, MMI Server, WebApp), com formatos diferentes:
\begin{itemize}
    \item Java requer keystores JKS
    \item Python requer certificados PEM
\end{itemize}

\textbf{Solução}:
\begin{enumerate}
    \item Geração de certificados self-signed usando OpenSSL
    \item Conversão para formatos apropriados:
    \begin{verbatim}
openssl req -x509 -newkey rsa:4096 -keyout key.pem -out cert.pem -days 365
keytool -importcert -file cert.pem -keystore keystore2.jks
    \end{verbatim}
    \item Cópia de certificados para diretórios corretos de cada componente
    \item Configuração de SSL contexts para ignorar verificação (ambiente de desenvolvimento)
\end{enumerate}

\textbf{Lição}: Planeamento de PKI (Public Key Infrastructure) desde início do projeto.

\subsection{Sincronização de Componentes Distribuídos}

\textbf{Problema}: Iniciar 5 processos independentes (FUSION, MMI, RASA, WEBAPP, SERVER) manualmente é propenso a erros de ordem e timing.

\textbf{Solução}: Criação de script \texttt{start.bat} que:
\begin{enumerate}
    \item Lança todos os componentes em Windows Terminals separados
    \item Utiliza caminhos relativos para portabilidade
    \item Ativa ambiente virtual automaticamente
    \item Abre browser com interface web após componentes iniciarem
\end{enumerate}

\textbf{Lição}: Automação de deployment é essencial para sistemas distribuídos, mesmo em ambiente de desenvolvimento.

\subsection{Mudanças Frequentes na Interface do Google Maps}

\textbf{Problema}: Google Maps atualiza estrutura HTML/CSS frequentemente, quebrando seletores Selenium.

\textbf{Solução}:
\begin{enumerate}
    \item Uso de seletores robustos (ARIA labels, data attributes) em vez de classes CSS
    \item Implementação de múltiplos seletores fallback:
    \begin{verbatim}
SEARCH_BOX_SELECTORS = [
    (By.ID, "searchboxinput"),
    (By.CSS_SELECTOR, "input[aria-label*='Search']"),
    (By.XPATH, "//input[@name='q']")
]
    \end{verbatim}
    \item Retry logic que tenta diferentes estratégias de localização
    \item Logging detalhado de elementos não encontrados para debugging rápido
\end{enumerate}

\textbf{Lição}: Web scraping/automation requer manutenção contínua. Considerar APIs oficiais quando disponíveis.

\subsection{TTS a Falar Duas Vezes}

\textbf{Problema}: Mensagens TTS eram reproduzidas duas vezes: uma vez no Assistant e outra no browser.

\textbf{Causa}: WebSocket server estava broadcasting mensagem de volta ao sender.

\textbf{Solução}: Filtrar sender ao fazer broadcast:
\begin{verbatim}
for client in connected_clients:
    if client != websocket:  # Don't echo back
        await client.send(message)
\end{verbatim}

\textbf{Lição}: Testar fluxos de comunicação bidirecional cuidadosamente para evitar loops.

\subsection{Resolução de Dependências RASA}

\textbf{Problema}: Instalação de RASA em ambiente virtual provocava resolução de dependências extremamente lenta (> 30 minutos) devido a conflitos entre centenas de pacotes.

\textbf{Solução}: Utilizar instalação global de RASA existente em Python 3.10, instalando apenas dependências essenciais no venv (selenium, websockets, colorlog).

\textbf{Lição}: Para frameworks pesados de ML, considerar instalações globais ou uso de containers Docker para isolamento.

\begin{comment}
\section{Trabalho Futuro}

Diversas melhorias e extensões podem ser implementadas para expandir as capacidades do sistema:

\subsection{Melhorias na Fusão}

\begin{itemize}
    \item \textbf{Aprendizagem de Preferências}: Sistema que aprende padrões de uso do utilizador e adapta regras de fusão (e.g., utilizador prefere sempre voz para pesquisas)

    \item \textbf{Fusão Baseada em Machine Learning}: Substituir regras hard-coded por modelo neural que aprende a fundir modalidades baseado em dados de uso

    \item \textbf{Janela Temporal Adaptativa}: Ajustar duração da janela de fusão baseado no tipo de comando e contexto (comandos complexos podem requerer janelas maiores)

    \item \textbf{Fusão de 3+ Modalidades}: Adicionar eye-tracking para identificar pontos de interesse no mapa, ou toque em tablet para especificar localizações

    \item \textbf{Contexto Conversacional}: Manter histórico de conversas para resolver referências anafóricas (e.g., "mostra-me restaurantes" seguido de "quero avaliações disso")
\end{itemize}

\subsection{Expansão de Modalidades}

\begin{itemize}
    \item \textbf{Reconhecimento de Emoções}: Análise de cadência na voz para detetar frustração/urgência e adaptar comportamento

    \item \textbf{Gestos Dinâmicos Complexos}: Reconhecimento de gestos contínuos (e.g., traçar rota com mão) usando redes recorrentes (LSTM/GRU)
    \item \textbf{Toque Multimodal}: Integração com touchscreen para permitir fusão toque + voz (e.g., "zoom aqui" + toque no ecrã)

    \item \textbf{Gaze Tracking}: Eye-tracking para identificar automaticamente pontos de interesse sem necessidade de apontar

    \item \textbf{Deteção de Atividade}: Usar acelerómetro/giroscópio para detetar se utilizador está a pé, em veículo, etc., e adaptar interface
\end{itemize}

\subsection{Melhorias na Interface e Usabilidade}

\begin{itemize}
    \item \textbf{Visualização de Estado de Fusão}: Interface gráfica que mostra em tempo real quais modalidades estão ativas, eventos recebidos, e decisões de fusão

    \item \textbf{Modo de Treino Personalizado}: Sistema que permite utilizador definir gestos customizados através de demonstração

    \item \textbf{Dashboard de Analytics}: Painel que mostra estatísticas de uso (modalidades mais usadas, comandos frequentes, taxa de sucesso)

    \item \textbf{Configuração Gráfica}: GUI para ajustar parâmetros do sistema (janela temporal, thresholds de confidence, priorização de modalidades) sem editar código

    \item \textbf{Multi-Utilizador}: Suporte para múltiplos utilizadores simultâneos com perfis individuais de preferências
\end{itemize}

\subsection{Robustez e Performance}

\begin{itemize}
    \item \textbf{Caching Inteligente}: Cache de resultados de pesquisa e direções para reduzir latência em queries repetidas

    \item \textbf{Recuperação Sem Reinício}: Sistema de checkpointing que permite recuperar de falhas sem reiniciar todos os componentes

    \item \textbf{Otimização de Pipeline de Gestos}: Reduzir latência de reconhecimento de gestos através de:
    \begin{itemize}
        \item Processamento em GPU (CUDA)
        \item Algoritmos de early classification (decidir antes de gesto terminar)
        \item Modelos mais leves (MobileNet para deteção de poses)
    \end{itemize}

    \item \textbf{Testes End-to-End Automatizados}: Suite de testes que simula comandos multimodais e valida execução correta

    \item \textbf{Deployment em Container}: Dockerização de componentes para deployment simplificado e escalabilidade
\end{itemize}
    
\section{Aprendizagens}

Este projeto proporcionou conhecimentos valiosos em múltiplas áreas:

\subsection{Pensamento Crítico}

\begin{itemize}
    \item Como atribuir valor real ao uso intermodal de sistemas
    \item Situações reais em que este não seria uma desvantagem
    \item Gestos e comandos perceptíveis, e que cumpram os requisitos estudados
\end{itemize}

\subsection{Arquitetura de Sistemas Multimodais}

\begin{itemize}
    \item Design de arquiteturas distribuídas com comunicação assíncrona
    \item Trade-offs entre fusão centralizada vs. distribuída
    \item Gestão de estado em sistemas event-driven
\end{itemize}

\subsection{Processamento de Linguagem Natural}

\begin{itemize}
    \item Treino e fine-tuning de modelos RASA para domínios específicos
    \item Importância de corpus de treino diversificado e bem anotado
    \item Gestão de confiança e ambiguidade em reconhecimento de intents
\end{itemize}

\subsection{Reconhecimento de Gestos}

\begin{itemize}
    \item Processamento de dados skeleton 3D do Kinect
    \item Algoritmos de segmentação temporal de gestos
    \item Importância de thresholds adaptativos baseados em contexto
\end{itemize}

\subsection{Fusão de Informação}

\begin{itemize}
    \item Estratégias de fusão temporal vs. semântica
    \item Resolução de conflitos em sistemas multimodais
    \item Validação e verificação de comandos fundidos
\end{itemize}

\subsection{Automação Web}

\begin{itemize}
    \item Técnicas avançadas de Selenium WebDriver
    \item Estratégias para lidar com SPAs (Single Page Applications) dinâmicas
    \item Robustez através de retry logic e seletores múltiplos
    \item Execução de JavaScript para controlo preciso
\end{itemize}

\subsection{Engenharia de Software}

\begin{itemize}
    \item Importância de logging detalhado para debugging de sistemas distribuídos
    \item Gestão de dependências e compatibilidade de versões
    \item Automação de deployment e inicialização de sistemas complexos
    \item Documentação como parte integral do desenvolvimento
\end{itemize}
\end{comment}

\section{Reflexão Final}

O sistema desenvolvido demonstra a viabilidade e os benefícios de interfaces multimodais para controlo de aplicações complexas. A combinação de voz e gestos proporciona uma experiência de utilizador significativamente mais natural e eficiente do que interfaces tradicionais baseadas exclusivamente em rato e teclado.

A arquitetura modular baseada em componentes independentes e protocolo MMI padronizado facilita a extensão futura do sistema com novas modalidades ou funcionalidades. O código desenvolvido segue boas práticas de engenharia de software, com separação clara de responsabilidades, tratamento robusto de erros, e documentação inline.

Este trabalho estabelece uma base sólida para investigação futura em interação multimodal, podendo ser expandido para outros domínios de aplicação além de mapas digitais. O sistema desenvolvido tem potencial para ser utilizado em contextos reais, nomeadamente em ambientes onde interação hands-free é vantajosa (e.g., condução, acessibilidade).

A crescente ubiquidade de sensores (câmaras, microfones, acelerómetros) em dispositivos modernos torna interfaces multimodais cada vez mais relevantes. Este projeto contribui para a compreensão de como combinar eficazmente múltiplas modalidades para criar experiências de utilizador ricas e naturais.


%%%%%%%%%%%%%%%%%%%%%%%%%%%%%%%%%
\printbibliography

\end{document}